% Section file for Chapter: Twisted Holography and Quantum Gravity
% Result (G9): Critical String Dichotomy

\section{Critical string dichotomy}
% label removed: sec:thqg-critical-string-dichotomy
\label{ch:thqg-critical-string}% alias for dangling \ref{ch:...}
\index{critical string!dichotomy|textbf}
\index{Virasoro!c=26@$c=26$ critical string}
\index{transgression algebra|textbf}
\index{Clifford completion|textbf}

The Virasoro algebra at $c = 26$ has appeared in three earlier contexts:
as the uncurved dual $\mathrm{Vir}_{26}^! \simeq \mathrm{Vir}_0$
(Proposition~\ref{V1-prop:virasoro-c26-selfdual}),
as the endpoint of the complementarity sum
$\kappa + \kappa' = c/2 + (26-c)/2 = 13$
(Proposition~\ref{V1-prop:vir-complementarity}),
and as the saturation point of the depth-zero resonance shadow
(Remark~\ref{V1-rem:virasoro-resonance-model}).
The three facts coalesce into a single structural dichotomy,
controlled by a Clifford algebra built from the handle
decomposition of genus-$g$ surfaces.

\begin{remark}[Scope restriction inherited from Volume~I]
The \ClaimStatusProvedHere{} results of this section build on Volume~I
statements (Theorems \texttt{vir-all-genera}, \texttt{topological-regime},
\texttt{gravitational-regime}, \texttt{critical-string-dichotomy},
\texttt{ghost-sector-c26}, \texttt{vir-self-duality-c13};
Propositions \texttt{vir-complementarity},
\texttt{virasoro-c26-selfdual},
\texttt{genus-g-curvature-package}) whose scope qualifiers propagate
to the present proofs. Load-bearing cross-volume citations in proof
bodies below carry the explicit ``Volume~I'' prose prefix;
the semantic content (complementarity tallies, regime dichotomy,
secondary-anomaly data) is inherited without restatement.
\end{remark}

The construction proceeds from the curved bar complex alone.
The curvature $\dfib^{\,2} = \kappa(\cA) \cdot \omega_g$
(Convention~\ref{V1-conv:higher-genus-differentials})
is not an obstruction to be cancelled but a datum to be \emph{transgressed}:
the transgression algebra $B_\Theta$ is the universal solution to the
problem of killing the curvature by adjoining an element whose square
\emph{is} the curvature. The secondary anomaly element
$u = \eta^2 \in B_\Theta$ is closed and central. Its invertibility or
vanishing controls everything: when $u$ is invertible, the genus-$g$
Clifford completion is Morita equivalent to $B_\Theta[u^{-1}]$ via the
spinor module; when $u = 0$, the Clifford relations degenerate to
exterior algebra relations, and the bar complex becomes a genuine cochain
complex tensored with $\bigwedge^{2g+1}$. For the Virasoro at $c = 0$
(i.e., the Koszul dual of $\mathrm{Vir}_{26}$),
$u = 0$ is the algebraic exterior/uncurved shadow of the ghost-sector
simplification that appears in the no-ghost theorem;
at $c = 13$, the symmetric distribution $\kappa = \kappa' = 13/2$ gives
the self-dual fixed point of the Feigin--Frenkel involution.

% ======================================================================
%
% SUBSECTION: THE TRANSGRESSION ALGEBRA
%
% ======================================================================

\subsection{The transgression algebra}
% label removed: subsec:transgression-algebra
\index{transgression algebra!construction|textbf}

\subsubsection{Construction from the curved bar complex}
% label removed: subsubsec:transgression-construction

The genus-$g$ bar complex $(\barB^{(g)}(\cA), \dfib)$ of a cyclic
chiral algebra~$\cA$ satisfies
$\dfib^{\,2} = \kappa(\cA) \cdot \omega_g \cdot \id$
(Proposition~\ref{V1-prop:genus-g-curvature-package}),
with $\kappa(\cA) = \mathrm{Tr}_\cA$
(Definition~\ref{V1-def:modular-curvature-direct}) and $\omega_g$ the
Arakelov $(1,1)$-form on the genus-$g$ fiber.
The transgression algebra is the universal solution to the
problem of killing this curvature by adjoining a degree-one
element whose square is the curvature.

\begin{definition}[Transgression algebra]
% label removed: def:transgression-algebra
\index{transgression algebra!definition|textbf}
Let $(B, d)$ be a graded-associative $\Bbbk$-algebra
equipped with a degree-$1$ derivation $d\colon B \to B$ satisfying the
\emph{curved} relation
\begin{equation}% label removed: eq:g9-curved-relation
d^2(b) = [\lambda, b]
\qquad
\text{for all } b \in B,
\end{equation}
where $\lambda \in B^2$ is a degree-$2$ central element (the
\emph{curvature}). The \emph{transgression algebra} of the pair
$(B, d, \lambda)$ is
\begin{equation}% label removed: eq:g9-transgression-def
B_\Theta
\;:=\;
B\langle \eta \rangle \big/ (\eta^2 - \lambda),
\end{equation}
where $\eta$ is a formal generator of degree~$1$, subject to
$|\eta| = 1$ and $\eta^2 = \lambda$. The differential extends to
$B_\Theta$ by
\begin{equation}% label removed: eq:g9-transgression-diff
d_{B_\Theta}(\eta) := 0,
\qquad
d_{B_\Theta}(b) := d(b) + [\eta, b]
\quad\text{for } b \in B.
\end{equation}
\end{definition}

\begin{remark}[Degree and sign conventions]
% label removed: rem:g9-transgression-signs
In~\eqref{V1-eq:g9-transgression-def}, the algebra $B\langle \eta \rangle$
is the free graded-associative extension of~$B$ by an odd generator
$\eta$ of degree~$1$, with the Koszul sign rule: $\eta \cdot b =
(-1)^{|b|} b \cdot \eta$ for homogeneous $b \in B$. The relation
$\eta^2 = \lambda$ is consistent with graded commutativity: since
$|\eta| = 1$ is odd, $\eta^2 = \eta \cdot \eta$ is
\emph{not} forced to vanish in the graded-commutative world (where it
would be $-\eta^2$); in the associative setting, $\eta^2$ is a
well-defined degree-$2$ element. The relation $\eta^2 = \lambda$
identifies this element with the curvature.
\end{remark}

\begin{proposition}[Transgression kills curvature; \ClaimStatusProvedHere]
% label removed: prop:transgression-kills-curvature
\index{transgression algebra!kills curvature}
The extended differential $d_{B_\Theta}$ on $B_\Theta$ satisfies
\begin{equation}% label removed: eq:g9-dBT-squared
d_{B_\Theta}^{\,2} = 0.
\end{equation}
\end{proposition}

\begin{proof}
We verify $d_{B_\Theta}^{\,2}$ on generators.

\emph{On $\eta$}:
$d_{B_\Theta}(\eta) = 0$, so $d_{B_\Theta}^{\,2}(\eta) = 0$.

\emph{On $b \in B$}: Compute
\begin{align}
d_{B_\Theta}^{\,2}(b)
&= d_{B_\Theta}\bigl(d(b) + [\eta, b]\bigr) \notag \\
&= d_{B_\Theta}(d(b)) + d_{B_\Theta}([\eta, b]) \notag \\
&= \bigl(d^2(b) + [\eta, d(b)]\bigr)
 + \bigl([d_{B_\Theta}(\eta), b]
 + (-1)^{|\eta|}[\eta, d_{B_\Theta}(b)]\bigr). % label removed: eq:g9-dsq-expand
\end{align}
The first term in the second parenthesis vanishes since
$d_{B_\Theta}(\eta) = 0$. Expand the remaining terms:
\begin{align}
d_{B_\Theta}^{\,2}(b)
&= d^2(b) + [\eta, d(b)]
 - [\eta, d(b) + [\eta, b]] \notag \\
&= d^2(b) + [\eta, d(b)]
 - [\eta, d(b)] - [\eta, [\eta, b]] \notag \\
&= d^2(b) - [\eta, [\eta, b]]. % label removed: eq:g9-dsq-simplify
\end{align}
The sign $(-1)^{|\eta|} = -1$ accounts for the negative sign in the
last line. Now compute the double bracket.
Using associativity and the relation $\eta^2 = \lambda$:
\begin{align}
[\eta, [\eta, b]]
&= \eta(\eta b - (-1)^{|b|}b\eta)
 - (-1)^{|b|+1}(\eta b - (-1)^{|b|}b\eta)\eta \notag \\
&= \eta^2 b - (-1)^{|b|}\eta b \eta
 - (-1)^{|b|+1}\eta b \eta + (-1)^{2|b|+1}b\eta^2 \notag \\
&= \lambda b - (-1)^{|b|}\eta b \eta
 + (-1)^{|b|}\eta b \eta - b\lambda \notag \\
&= \lambda b - b \lambda \notag \\
&= [\lambda, b]. % label removed: eq:g9-bracket-curvature
\end{align}
Substituting into~\eqref{V1-eq:g9-dsq-simplify}:
\[
d_{B_\Theta}^{\,2}(b)
= d^2(b) - [\lambda, b]
= [\lambda, b] - [\lambda, b]
= 0,
\]
using $d^2(b) = [\lambda, b]$ from~\eqref{V1-eq:g9-curved-relation}.
Since $d_{B_\Theta}$ is a derivation and vanishes on all
generators, $d_{B_\Theta}^{\,2} = 0$ on all of $B_\Theta$.
\end{proof}

\begin{remark}[Module structure and $\eta$-decomposition]
% label removed: rem:g9-eta-filtration
\index{transgression algebra!module structure}
Since $\eta^2 = \lambda \in B$, any element of $B_\Theta$ can be
written uniquely as $b_0 + b_1 \eta$ with $b_0, b_1 \in B$.
Thus $B_\Theta$ is a rank-$2$ free left $B$-module:
\begin{equation}% label removed: eq:g9-BT-decomposition
B_\Theta = B \oplus B\eta.
\end{equation}
The multiplication rule is
\begin{equation}% label removed: eq:g9-BT-multiplication
(a_0 + a_1\eta)(b_0 + b_1\eta)
= (a_0 b_0 + (-1)^{|a_1|}a_1 b_1 \lambda)
 + (a_0 b_1 + (-1)^{|a_1|}a_1 b_0)\eta.
\end{equation}
\end{remark}

\begin{proposition}[Quotient identification; \ClaimStatusProvedHere]
% label removed: prop:g9-quotient
\index{transgression algebra!quotient}
The quotient by the ideal generated by~$\eta$ recovers the algebra
modulo its curvature:
\begin{equation}% label removed: eq:g9-quotient-id
B_\Theta / (\eta) \;\cong\; B / (\lambda).
\end{equation}
The induced differential on $B/(\lambda)$ squares to zero.
\end{proposition}

\begin{proof}
From~\eqref{V1-eq:g9-BT-decomposition}, the ideal $(\eta) = B\eta + B\eta^2
= B\eta + B\lambda$. The quotient
$B_\Theta/(\eta) = (B \oplus B\eta)/(B\eta + B\lambda)
\cong B/(B\lambda) = B/(\lambda)$.
On the quotient, $d_{B_\Theta}(b) = d(b) + [\eta, b] \mapsto d(b)$
(since $\eta \mapsto 0$), and $d^2(b) = [\lambda, b] \mapsto 0$ since
$\lambda \mapsto 0$ in $B/(\lambda)$.
\end{proof}

\subsubsection{Specialization to the bar complex}
% label removed: subsubsec:transgression-bar
\index{transgression algebra!bar complex specialization}

Apply the transgression construction to the genus-$g$ bar complex of a
cyclic chiral algebra~$\cA$. Set:
\begin{equation}% label removed: eq:g9-specialization
B := \barB^{(g)}(\cA),
\quad
d := \dfib,
\quad
\lambda := \kappa(\cA) \cdot \omega_g.
\end{equation}
The curvature $\lambda = \kappa \cdot \omega_g$ is central: it is a scalar
multiple of the identity in the bar coalgebra tensored with the
$(1,1)$-class on the fiber.

\begin{definition}[Bar-transgression complex]
% label removed: def:bar-transgression-complex
\index{bar-transgression complex|textbf}
The \emph{bar-transgression complex} of $\cA$ at genus $g$ is
\begin{equation}% label removed: eq:g9-bar-transgression
B_\Theta^{(g)}(\cA)
\;:=\;
\barB^{(g)}(\cA)\langle \eta \rangle
\big/ (\eta^2 - \kappa(\cA) \cdot \omega_g),
\end{equation}
with $d_{B_\Theta}^{\,2} = 0$ by
Proposition~\ref{V1-prop:transgression-kills-curvature}.
Its cohomology is the \emph{transgressed bar cohomology}
$H^*(B_\Theta^{(g)}(\cA))$.
\end{definition}

\begin{remark}[Relation to the MC element]
% label removed: rem:g9-mc-relation
\index{transgression algebra!MC element relation}
The curvature $\kappa(\cA) \cdot \omega_g$ is the genus-$g$ scalar
shadow of the universal MC element $\Theta_\cA$:
$\kappa(\cA) \cdot \omega_g = \mathrm{Sh}_{g,0}(\Theta_\cA)$
(Corollary~\ref*{V1-cor:shadow-extraction}).
The transgression algebra is the universal algebraic solution to
killing this shadow. The corrected differential $\Dg{g}$
(Convention~\ref{V1-conv:higher-genus-differentials}) is a
\emph{geometric} solution using $A$-cycle periods. The universal
property ensures a comparison map
\[
B_\Theta^{(g)}(\cA) \longrightarrow
(\barB^{(g)}(\cA),\, \Dg{g}),
\qquad
\eta \longmapsto
\sum_{k=1}^g t_k\, \omega_k,
\]
which is a dg algebra homomorphism when the $A$-periods are specified.
\end{remark}

% ======================================================================
%
% SUBSECTION: THE SECONDARY ANOMALY ELEMENT
%
% ======================================================================

\subsection{The secondary anomaly element}
% label removed: subsec:secondary-anomaly
\index{secondary anomaly!construction|textbf}

\begin{definition}[Secondary anomaly element]
% label removed: def:secondary-anomaly
\index{secondary anomaly!definition|textbf}
The \emph{secondary anomaly element} of the transgression algebra
$B_\Theta^{(g)}(\cA)$ is
\begin{equation}% label removed: eq:g9-secondary-anomaly
u \;:=\; \eta^2 \;=\; \kappa(\cA) \cdot \omega_g
\;\in\; B_\Theta^{(g)}(\cA).
\end{equation}
It has cohomological degree $|u| = 2$.
\end{definition}

\begin{proposition}[Properties of $u$; \ClaimStatusProvedHere]
% label removed: prop:secondary-anomaly-properties
\index{secondary anomaly!properties}
The secondary anomaly element satisfies:
\begin{enumerate}[label=\textup{(\roman*)}]
\item \emph{Closedness.}\; $d_{B_\Theta}(u) = 0$.

\item \emph{Centrality.}\; $u \cdot b = b \cdot u$ for all $b \in
B_\Theta^{(g)}(\cA)$.

\item \emph{Factorization.}\; $u = \kappa(\cA) \cdot \omega_g$, where
$\kappa(\cA) \in \Bbbk$ is the scalar modular characteristic and
$\omega_g \in H^{1,1}(\Sigma_g)$ is the Arakelov class. The element
$u$ vanishes if and only if $\kappa(\cA) = 0$.

\item \emph{Shadow obstruction tower level.}\; Under the identification
$\kappa(\cA) = \Theta_\cA^{\leq 2}$
(Definition~\ref{V1-def:modular-curvature-direct}), the
secondary anomaly is the degree-$2$ shadow evaluated on the genus-$g$
fiber: $u = \Theta_\cA^{\leq 2}\big|_{\Sigma_g}$.
\end{enumerate}
\end{proposition}

\begin{proof}
\emph{(i)} Compute $d_{B_\Theta}(u) = d_{B_\Theta}(\eta^2)$. By the
Leibniz rule,
\[
d_{B_\Theta}(\eta^2) = d_{B_\Theta}(\eta) \cdot \eta
+ (-1)^{|\eta|}\eta \cdot d_{B_\Theta}(\eta) = 0 + 0 = 0,
\]
since $d_{B_\Theta}(\eta) = 0$.

\emph{(ii)} For $b \in B$: since $u = \kappa \cdot \omega_g$ is a
scalar in the center of~$B$, we have $u \cdot b = b \cdot u$. For
$b = \eta$: $u \cdot \eta = \eta^2 \cdot \eta = \eta \cdot \eta^2 =
\eta \cdot u$, where the middle step uses centrality of
$\kappa \cdot \omega_g$ in $B$. Since $B_\Theta = B \oplus B\eta$,
$u$ commutes with everything.

\emph{(iii)} The factorization is the definition. The Arakelov form
$\omega_g$ is strictly positive on $\Sigma_g$ for $g \geq 1$
(Proposition~\ref{V1-prop:genus-g-curvature-package}), so $u = 0$
iff $\kappa(\cA) = 0$.

\emph{(iv)} By Definition~\ref{V1-def:modular-curvature-direct},
$\kappa(\cA) = \Theta_\cA^{\leq 2}$ is the weight-$2$ shadow.
\end{proof}

\begin{remark}[The element $u$ as characteristic class]
% label removed: rem:g9-u-characteristic
\index{secondary anomaly!as characteristic class}
Over the moduli space $\overline{\cM}_g$, the element $u$ defines a
section of the Hodge line bundle scaled by $\kappa$:
\[
u \;\in\; H^0\bigl(\overline{\cM}_g,\;
\kappa(\cA) \cdot c_1(\mathbb{E})\bigr).
\]
Its vanishing $\{u = 0\}$ is the
\emph{uncurved locus} $\{\kappa(\cA) = 0\}$: the algebraic
content of anomaly cancellation.
\end{remark}

% ======================================================================
%
% SUBSECTION: CONNECTING HOMOMORPHISM
%
% ======================================================================

\subsection{The secondary anomaly as connecting homomorphism}
% label removed: subsec:secondary-connecting
\index{secondary anomaly!connecting homomorphism}

The transgression algebra fits into a short exact sequence whose
connecting homomorphism is a Bockstein operator.

\begin{proposition}[Short exact sequence of the transgression;
\ClaimStatusProvedHere]
% label removed: prop:g9-ses-transgression
There is a short exact sequence of cochain complexes
\begin{equation}% label removed: eq:g9-ses
0 \longrightarrow
B \xrightarrow{\;\cdot\,\eta\;}
B_\Theta
\xrightarrow{\;\mathrm{pr}\;}
B / (\lambda)
\longrightarrow 0,
\end{equation}
where the first map is multiplication by $\eta$ and the second is
projection modulo the ideal $(\eta)$. The connecting homomorphism
in the associated long exact sequence is
\begin{equation}% label removed: eq:g9-connecting
\delta\colon H^n(B/(\lambda)) \longrightarrow H^{n+1}(B),
\qquad
\delta([\bar{b}]) = [d(b)],
\end{equation}
where $b \in B$ is any lift of $\bar{b} \in B/(\lambda)$.
\end{proposition}

\begin{proof}
Exactness at each term follows from the rank-$2$ decomposition
$B_\Theta = B \oplus B\eta$: the image of $\cdot\eta$ is $B\eta$,
which is the kernel of $\mathrm{pr}$. Surjectivity of $\mathrm{pr}$
and injectivity of $\cdot\eta$ (since $B$ is torsion-free over
$\Bbbk$) give the short exact sequence.

For the connecting homomorphism: let $\bar{b} \in B/(\lambda)$ be a
cocycle. Lift to $b \in B$. Then
$d_{B_\Theta}(b) = d(b) + [\eta, b] \in B_\Theta$. The
$B\eta$-component of this is $[\eta, b]$, so the projection to
$B/(\lambda)$ sends $d_{B_\Theta}(b) \mapsto d(b)$ mod~$(\lambda)$.
The closedness condition means $d(b) \equiv 0 \pmod{\lambda}$, so
$d(b) = \lambda \cdot c$ for some $c \in B$. The connecting
homomorphism takes $[\bar{b}]$ to $[d(b)]$ in $H^{n+1}(B)$.
\end{proof}

\begin{corollary}[Bockstein long exact sequence; \ClaimStatusProvedHere]
% label removed: cor:g9-bockstein
\index{Bockstein homomorphism!transgression}
The long exact sequence reads
\begin{equation}% label removed: eq:g9-bockstein-les
\cdots \to
H^n(B) \xrightarrow{\cdot\eta}
H^n(B_\Theta) \xrightarrow{\mathrm{pr}}
H^n(B/(\lambda)) \xrightarrow{\delta}
H^{n+1}(B) \to \cdots
\end{equation}
Classes in $H^*(B/(\lambda))$ that lift to the transgressed bar
cohomology $H^*(B_\Theta)$ are precisely those annihilated
by the Bockstein~$\delta$.
\end{corollary}

\begin{remark}[Spectral sequence of the $\eta$-filtration]
% label removed: rem:g9-eta-spectral
\index{spectral sequence!$\eta$-filtration}
The two-step filtration $F^0 B_\Theta = B_\Theta$,
$F^1 B_\Theta = B\eta$ gives a spectral sequence with
$E_1^{0,q} = H^q(B)$, $E_1^{1,q} = H^{q-1}(B)$, and
$d_1$ is multiplication by $u = \kappa \cdot \omega_g$:
\[
E_2^{0,q} = \ker(u\colon H^q(B) \to H^{q+2}(B)),
\qquad
E_2^{1,q} = \mathrm{coker}(u\colon H^{q-1}(B) \to H^{q+1}(B)).
\]
Degeneration at $E_2$ (forced by the two-step filtration) is the
algebraic skeleton of the two-regime dichotomy.
\end{remark}

% ======================================================================
%
% SUBSECTION: THE CLIFFORD COMPLETION
%
% ======================================================================

\subsection{The Clifford completion}
% label removed: subsec:clifford-completion
\index{Clifford completion!construction|textbf}
\index{handle operators|textbf}

The genus-$g$ surface $\Sigma_g$ has first homology
$H_1(\Sigma_g, \Z) \cong \Z^{2g}$ with the intersection pairing
$\langle -, - \rangle\colon H_1 \otimes H_1 \to \Z$.
There exists a symplectic basis $\{A_1, B_1, \ldots, A_g, B_g\}$ with
\begin{equation}% label removed: eq:g9-intersection
\langle A_i, B_j \rangle = \delta_{ij},
\qquad
\langle A_i, A_j \rangle = 0,
\qquad
\langle B_i, B_j \rangle = 0.
\end{equation}
The sewing construction of the genus-$g$ bar complex from genus-$0$
data passes through $g$ handle attachments, one for each pair
$(A_k, B_k)$.

\begin{definition}[Handle operators]
% label removed: def:handle-operators
\index{handle operators!definition|textbf}
Let $B_\Theta = B_\Theta^{(g)}(\cA)$ be the bar-transgression complex
at genus~$g$. For each $k = 1, \ldots, g$, define the
\emph{handle operators}
\begin{equation}% label removed: eq:g9-handle-ops
\alpha_k \;:=\; \oint_{A_k} \eta^{(g)},
\qquad
\beta_k \;:=\; \oint_{B_k} \eta^{(g)},
\end{equation}
where $\eta^{(g)}$ is the degree-$1$ transgression element of
$B_\Theta^{(g)}(\cA)$, and the integrals are fiber integrals along
the cycles $A_k$, $B_k \subset \Sigma_g$. Each handle operator has
degree~$1$ and acts on $B_\Theta^{(g)}(\cA)$ by left multiplication.
\end{definition}

\begin{remark}[Geometric origin of handle operators]
% label removed: rem:g9-handle-geometric
\index{handle operators!geometric origin}
The fiber integral $\oint_{A_k} \eta^{(g)}$ is the residue of the
transgression element along the $A_k$-cycle. In the sewing
construction (Theorem~\ref{thm:heisenberg-sewing}), each handle
attachment corresponds to a self-sewing of the genus-$0$ surface:
cutting along a pair of cycles and regluing. The handle operator
$\alpha_k$ records the $A$-cycle contraction and $\beta_k$ the
$B$-cycle contraction.
\end{remark}

\begin{theorem}[Clifford relations for handle operators;
\ClaimStatusProvedHere]
% label removed: thm:clifford-relations
\index{Clifford algebra!handle operators|textbf}
The handle operators satisfy:
\begin{align}
\alpha_i \beta_j + \beta_j \alpha_i &= \delta_{ij}\, u,
% label removed: eq:g9-clifford-ab \\
\alpha_i \alpha_j + \alpha_j \alpha_i &= 0,
% label removed: eq:g9-clifford-aa \\
\beta_i \beta_j + \beta_j \beta_i &= 0,
% label removed: eq:g9-clifford-bb
\end{align}
where $u = \kappa(\cA) \cdot \omega_g$ is the secondary anomaly
element. The generators
$\{\alpha_1, \beta_1, \ldots, \alpha_g, \beta_g\}$
generate a Clifford algebra $\mathrm{Cl}_{2g}(u)$ over the ring
$\Bbbk[u]$.
\end{theorem}

\begin{proof}
The proof reduces to the intersection pairing on $H_1(\Sigma_g, \Z)$
via the sewing calculus.

\emph{Step~1 (Mixed anticommutator, $i \neq j$).}\;
When $i \neq j$, the cycles $A_i$ and $B_j$ have intersection
number $\langle A_i, B_j \rangle = 0$. The contractions act on
disjoint tensor factors of the bar coalgebra. Since $\alpha_i$ and
$\beta_j$ are odd operators (degree~$1$) on disjoint data,
$\alpha_i \beta_j = -\beta_j \alpha_i$, giving
$\alpha_i\beta_j + \beta_j\alpha_i = 0$.

\emph{Step~2 (Mixed anticommutator, $i = j$).}\;
The cycles $A_i$ and $B_i$ have $\langle A_i, B_i \rangle = 1$.
The iterated fiber integral of $(\eta^{(g)})^2 = u =
\kappa \cdot \omega_g$ over the pair $(A_i, B_i)$ evaluates via
the intersection formula:
\begin{equation}% label removed: eq:g9-clifford-compute
\alpha_i \beta_i + \beta_i \alpha_i
= \oint_{A_i \times B_i}
 (\eta^{(g)})^2
= u \cdot \langle A_i, B_i \rangle
= u.
\end{equation}
The last equality uses the normalization~\eqref{V1-eq:g9-intersection}
of the symplectic basis. The factor~$u$ appears because the
integrand is $(\eta^{(g)})^2 = u$, and the fiber integral computes
the intersection number.

\emph{Step~3 (Same-type anticommutators).}\;
$\langle A_i, A_j \rangle = 0$ for all $i, j$, so
$\alpha_i \alpha_j + \alpha_j \alpha_i = 0$. Identically for
$\beta_i \beta_j$.

\emph{Step~4 (Clifford identification).}\;
Relations~\eqref{V1-eq:g9-clifford-ab}--\eqref{V1-eq:g9-clifford-bb}
are the defining relations of $\mathrm{Cl}_{2g}(q)$, where
$q\colon \Bbbk^{2g} \to \Bbbk$ is the quadratic form
$q(x) = u \cdot Q(x)$ with $Q$ the intersection form:
\[
J = \begin{pmatrix}
0 & I_g \\ -I_g & 0
\end{pmatrix}
\]
in the basis $(\alpha_1, \ldots, \alpha_g, \beta_1, \ldots, \beta_g)$.
The associated symmetric bilinear form
$\langle x, y \rangle_{\mathrm{Cl}} = u \cdot \delta_{ij}$ on the
standard basis follows from the symplectic-to-Clifford passage.
\end{proof}

\begin{definition}[Clifford completion]
% label removed: def:clifford-completion
\index{Clifford completion!definition|textbf}
The \emph{genus-$g$ Clifford completion} of $\cA$ is the algebra
\begin{equation}% label removed: eq:g9-clifford-completion
G_g(\cA)
\;:=\;
B_\Theta^{(g)}(\cA)
\;\otimes_{\Bbbk}\;
\mathrm{Cl}_{2g}(u),
\end{equation}
where $\mathrm{Cl}_{2g}(u)$ is the Clifford algebra generated by
$\{\alpha_1, \beta_1, \ldots, \alpha_g, \beta_g\}$ subject
to~\eqref{V1-eq:g9-clifford-ab}--\eqref{V1-eq:g9-clifford-bb}.
The differential extends by
$d_{G_g}(\alpha_k) = d_{G_g}(\beta_k) = 0$ and
$d_{G_g}(b) = d_{B_\Theta}(b)$ for $b \in B_\Theta^{(g)}(\cA)$.
Then $d_{G_g}^{\,2} = 0$.
\end{definition}

\begin{remark}[Dimension of the Clifford algebra]
% label removed: rem:g9-clifford-dimension
As a $\Bbbk$-vector space, $\mathrm{Cl}_{2g}(u)$ has dimension
$2^{2g}$. When $u$ is invertible,
$\mathrm{Cl}_{2g}(u) \cong \mathrm{Mat}_{2^g}(\Bbbk)$ by the
standard Morita classification. When $u = 0$,
$\mathrm{Cl}_{2g}(0) = \bigwedge^*(\Bbbk^{2g})$. This
$u = 0$ vs.\ $u \neq 0$ dichotomy is the structural divide.
\end{remark}

% ======================================================================
%
% SUBSECTION: EXPLICIT LOW-GENUS COMPUTATIONS
%
% ======================================================================

\subsection{Explicit computations at genus $1$ and $2$}
% label removed: subsec:g9-low-genus
\index{Clifford completion!low genus computations}

\subsubsection{Genus $1$: a single handle}
% label removed: subsubsec:g9-genus-1

At genus $1$, there is a single handle $(A_1, B_1)$ with
$\langle A_1, B_1 \rangle = 1$. The Clifford algebra is
$\mathrm{Cl}_2(u)$, generated by $\alpha := \alpha_1$,
$\beta := \beta_1$ with
\begin{equation}% label removed: eq:g9-genus1-clifford
\alpha\beta + \beta\alpha = u,
\qquad
\alpha^2 = 0,
\qquad
\beta^2 = 0.
\end{equation}

\emph{Multiplication table.}\;
As a $\Bbbk$-vector space, $\mathrm{Cl}_2(u)$ has basis
$\{1, \alpha, \beta, \alpha\beta\}$ with:
\begin{center}
\renewcommand{\arraystretch}{1.3}
\begin{tabular}{c|cccc}
$\cdot$ & $1$ & $\alpha$ & $\beta$ & $\alpha\beta$ \\
\hline
$1$ & $1$ & $\alpha$ & $\beta$ & $\alpha\beta$ \\
$\alpha$ & $\alpha$ & $0$ & $\alpha\beta$ & $0$ \\
$\beta$ & $\beta$ & $u - \alpha\beta$ & $0$ & $u\beta$ \\
$\alpha\beta$ & $\alpha\beta$ & $u\alpha$ & $0$ & $u\alpha\beta$
\end{tabular}
\end{center}

\emph{Volume element.}\;
$\mathrm{vol}_1 = \alpha\beta$ satisfies
$\mathrm{vol}_1^2 = \alpha\beta\alpha\beta =
\alpha(u - \alpha\beta)\beta = u\,\alpha\beta = u\,\mathrm{vol}_1$.

\emph{Spinor module.}\;
When $u$ is invertible, define the spinor module
$S_1 = \Bbbk \cdot v_+ \oplus \Bbbk \cdot v_-$ with
\begin{equation}% label removed: eq:g9-genus1-spinor
\alpha \cdot v_+ = v_-,
\quad
\alpha \cdot v_- = 0,
\quad
\beta \cdot v_+ = 0,
\quad
\beta \cdot v_- = u\, v_+.
\end{equation}
Verification: $(\alpha\beta + \beta\alpha) \cdot v_\pm =
u\, v_\pm$. Hence $S_1$ is a $\mathrm{Cl}_2(u)$-module, and
\begin{equation}% label removed: eq:g9-genus1-morita
\mathrm{Cl}_2(u) \cong \End_\Bbbk(S_1)
\cong \mathrm{Mat}_2(\Bbbk)
\qquad\text{(when $u$ is invertible).}
\end{equation}

\subsubsection{Genus $2$: two handles}
% label removed: subsubsec:g9-genus-2

At genus $2$, $H_1(\Sigma_2, \Z) \cong \Z^4$ has symplectic basis
$(A_1, B_1, A_2, B_2)$ with intersection matrix
\[
J_2 = \begin{pmatrix}
0 & 1 & 0 & 0 \\
-1 & 0 & 0 & 0 \\
0 & 0 & 0 & 1 \\
0 & 0 & -1 & 0
\end{pmatrix}.
\]
The Clifford algebra $\mathrm{Cl}_4(u)$ has
$\dim_\Bbbk = 2^4 = 16$, with basis
$\alpha_1^{a_1}\beta_1^{b_1}\alpha_2^{a_2}\beta_2^{b_2}$
for $a_i, b_i \in \{0,1\}$.

\emph{Spinor module.}\;
When $u$ is invertible, the spinor module $S_2 = \Bbbk^4$ has
basis $\{v_{++}, v_{+-}, v_{-+}, v_{--}\}$ with
\begin{alignat}{2}
\alpha_k \cdot v_{\ldots+_k\ldots}
 &= v_{\ldots-_k\ldots},
&\qquad
\alpha_k \cdot v_{\ldots-_k\ldots}
 &= 0, \notag \\
\beta_k \cdot v_{\ldots+_k\ldots}
 &= 0,
&\qquad
\beta_k \cdot v_{\ldots-_k\ldots}
 &= u\, v_{\ldots+_k\ldots}, % label removed: eq:g9-genus2-spinor
\end{alignat}
where the subscript indicates the $k$-th sign position. Hence
\begin{equation}% label removed: eq:g9-genus2-morita
\mathrm{Cl}_4(u) \cong \End_\Bbbk(S_2)
\cong \mathrm{Mat}_4(\Bbbk)
\qquad\text{(when $u$ is invertible).}
\end{equation}

\emph{Matrix units.}\;
The matrix units in terms of Clifford generators:
\begin{center}
\renewcommand{\arraystretch}{1.2}
\begin{tabular}{lll}
$E_{++,++} = u^{-2}\alpha_1\beta_1\alpha_2\beta_2$, &
$E_{++,--} = u^{-2}\alpha_1\alpha_2$, &
$E_{--,++} = \beta_1\beta_2$, \\
$E_{+-,-+} = u^{-1}\alpha_1\beta_2$, &
$E_{-+,+-} = u^{-1}\beta_1\alpha_2$, &
$E_{--,--} = u^{-2}\beta_1\alpha_1\beta_2\alpha_2$.
\end{tabular}
\end{center}
These $16$ matrix units span $\mathrm{Mat}_4$, confirming
$\mathrm{Cl}_4(u)[u^{-1}] \cong \mathrm{Mat}_4(\Bbbk[u^{\pm 1}])$.

% ======================================================================
%
% SUBSECTION: THE TWO REGIMES
%
% ======================================================================

\subsection{The two regimes: topological and gravitational}
% label removed: subsec:two-regimes
\index{critical string!two regimes|textbf}
\index{topological regime|textbf}
\index{gravitational regime|textbf}

The Clifford completion $G_g(\cA) = B_\Theta^{(g)}(\cA) \otimes
\mathrm{Cl}_{2g}(u)$ admits two fundamentally different structural
descriptions, determined by the invertibility of
$u = \kappa(\cA) \cdot \omega_g$.

\subsubsection{The topological regime: $u$ invertible}
% label removed: subsubsec:topological-regime

\begin{theorem}[Topological regime; \ClaimStatusProvedHere]
% label removed: thm:topological-regime
\index{topological regime!Morita equivalence|textbf}
Suppose $\kappa(\cA) \neq 0$, and that we work over a field
of characteristic~$0$ or characteristic not
dividing~$\kappa(\cA)$. Then:
\begin{enumerate}[label=\textup{(\roman*)}]
\item The Clifford completion is Morita equivalent to the
transgression algebra:
\begin{equation}% label removed: eq:g9-topo-morita
G_g(\cA)[u^{-1}]
\;\simeq_{\mathrm{Morita}}\;
\mathrm{Mat}_{2^g}\bigl(B_\Theta^{(g)}(\cA)[u^{-1}]\bigr).
\end{equation}

\item The Morita equivalence bimodule is the spinor module
\begin{equation}% label removed: eq:g9-spinor-bimod
S_g(\cA)
\;:=\;
B_\Theta^{(g)}(\cA)[u^{-1}] \otimes_\Bbbk S_g,
\end{equation}
where $S_g = \Bbbk^{2^g}$ is the standard spinor representation of
$\mathrm{Cl}_{2g}(u)$.

\item The derived categories are equivalent:
\begin{equation}% label removed: eq:g9-topo-derived
D(G_g(\cA)[u^{-1}]\text{-}\mathsf{mod})
\;\simeq\;
D(B_\Theta^{(g)}(\cA)[u^{-1}]\text{-}\mathsf{mod}).
\end{equation}
\end{enumerate}
\end{theorem}

\begin{proof}
\emph{(i)} The Clifford algebra $\mathrm{Cl}_{2g}(u)$ over a ring in
which $u$ is a unit is isomorphic to $\mathrm{Mat}_{2^g}$. Define the
canonical anticommutation operators
\begin{equation}% label removed: eq:g9-raising-lowering
a_k^+ := u^{-1}\alpha_k,
\qquad
a_k^- := \beta_k,
\qquad k = 1, \ldots, g.
\end{equation}
These satisfy $a_k^+ a_k^- + a_k^- a_k^+ = 1$,
$(a_k^+)^2 = (a_k^-)^2 = 0$, and $[a_k^\pm, a_l^\pm]_+ = 0$ for
$k \neq l$. The Fock space construction gives
$\mathrm{Cl}_{2g}(u)[u^{-1}] \cong \mathrm{Mat}_{2^g}(\Bbbk)$.

\emph{(ii)} The spinor module $S_g$ is the Fock space with vacuum
$|0\rangle := v_{+\cdots+}$ satisfying $\beta_k \cdot |0\rangle = 0$
for all $k$. The $2^g$ basis vectors are
$\alpha_{k_1}\cdots\alpha_{k_r}|0\rangle$ for
$1 \leq k_1 < \cdots < k_r \leq g$. Faithfulness of the left action
(since $u$ is a unit and the lowering operators have full range) gives
$\mathrm{Cl}_{2g}(u) \xrightarrow{\sim} \End(S_g)$.

\emph{(iii)} Morita-equivalent algebras have equivalent derived
categories; the functor $- \otimes_{G_g} S_g(\cA)$ provides it.
\end{proof}

\begin{remark}[Topological meaning of $2^g$]
% label removed: rem:g9-multiplicity
\index{topological regime!multiplicity}
The matrix size $2^g$ is the dimension of the half-spinor
representation, and $|H^1(\Sigma_g, \Z/2)| = 2^{2g}$ is the number of
spin structures on $\Sigma_g$. The Morita equivalence reduces the
Clifford completion to the transgression algebra at the cost of
choosing a spin structure (the vacuum $|0\rangle$).
\end{remark}

\subsubsection{The gravitational regime: $u = 0$}
% label removed: subsubsec:gravitational-regime

\begin{theorem}[Gravitational regime; \ClaimStatusProvedHere]
% label removed: thm:gravitational-regime
\index{gravitational regime!exterior algebra|textbf}
Suppose $\kappa(\cA) = 0$. Then:
\begin{enumerate}[label=\textup{(\roman*)}]
\item The transgression element satisfies $\eta^2 = 0$: the
bar complex is uncurved ($\dfib^{\,2} = 0$).

\item The Clifford relations degenerate to exterior algebra relations:
\[
\alpha_i \beta_j + \beta_j \alpha_i = 0,
\qquad
\alpha_i \alpha_j + \alpha_j \alpha_i = 0,
\qquad
\beta_i \beta_j + \beta_j \beta_i = 0.
\]

\item The Clifford completion splits:
\begin{equation}% label removed: eq:g9-grav-tensor
G_g(\cA)\big|_{u=0}
\;\cong\;
\barB^{(g)}(\cA)
\;\otimes\; \bigwedge\nolimits^*(\Bbbk^{2g}).
\end{equation}

\item The exterior algebra has rank $2^{2g}$ with top element
$\mathrm{vol}_g := \alpha_1 \beta_1 \cdots \alpha_g \beta_g
\in \bigwedge^{2g}(\Bbbk^{2g})$.
The Poincar\'e polynomial is $(1+t)^{2g}$.
\end{enumerate}
\end{theorem}

\begin{proof}
\emph{(i)} $u = \eta^2 = \kappa \cdot \omega_g = 0$.

\emph{(ii)} Set $u = 0$
in~\eqref{V1-eq:g9-clifford-ab}--\eqref{V1-eq:g9-clifford-bb}: all
anticommutators vanish.

\emph{(iii)} When $u = 0$ and $\lambda = 0$, the transgression
algebra quotient is $B_\Theta/(\eta) = B$ (uncurved), and
$\mathrm{Cl}_{2g}(0) = \bigwedge^*(\Bbbk^{2g})$.

\emph{(iv)} $\dim\bigwedge^k(\Bbbk^{2g}) = \binom{2g}{k}$;
$\sum_{k=0}^{2g}\binom{2g}{k}t^k = (1+t)^{2g}$.

The exterior algebra isomorphism in~(iii) is unique: over a
field, the exterior algebra on $2g$ generators is the unique
graded-commutative algebra with Hilbert series $(1+t)^{2g}$
(all generators in degree~$1$, all relations quadratic).
Any isomorphism $\mathrm{Cl}_{2g}(0) \cong
\bigwedge^*(\Bbbk^{2g})$ is determined by its restriction to
the degree-$1$ generators, where it is a linear isomorphism
$\Bbbk^{2g} \xrightarrow{\sim} \Bbbk^{2g}$; the algebra
structure forces the extension to all degrees.
\end{proof}

\begin{remark}[Ghost number in the gravitational regime]
% label removed: rem:g9-ghost-number
\index{ghost number!Clifford completion}
In the BRST language (Remark~\ref{V1-rem:bv-bar-bridge}), the generators
$\alpha_k$, $\beta_k$ carry ghost number~$1$. The volume element
$\mathrm{vol}_g$ has ghost number $2g$. In the gravitational regime,
the BRST cohomology splits by ghost number:
\[
H^*(G_g|_{u=0}) \cong
\bigoplus_{k=0}^{2g} H^{*-k}(\barB^{(g)}(\cA))
\otimes \bigwedge\nolimits^k(\Bbbk^{2g}).
\]
The ghost-number anomaly $2g = \dim H_1(\Sigma_g)$ counts the zero
modes of the ghost field. The exterior (rather than Clifford) structure
is equivalent to the vanishing of the ghost-number current anomaly.
\end{remark}

\begin{remark}[Explicit basis for $\bigwedge^*(\Bbbk^{2g})$]
% label removed: rem:g9-exterior-basis
\index{exterior algebra!explicit basis}
Basis elements are indexed by subsets $I \subset \{1, \ldots, 2g\}$:
for $I = \{i_1 < \cdots < i_k\}$, set $e_I = \gamma_{i_1}\cdots
\gamma_{i_k}$ where $(\gamma_1, \ldots, \gamma_{2g}) =
(\alpha_1, \beta_1, \ldots, \alpha_g, \beta_g)$.
Low-genus data:
\begin{center}
\renewcommand{\arraystretch}{1.2}
\begin{tabular}{ccl}
\toprule
$g$ & $\dim\bigwedge^*(\Bbbk^{2g})$ & Graded dimensions \\
\midrule
$1$ & $4$ & $1, 2, 1$ \\
$2$ & $16$ & $1, 4, 6, 4, 1$ \\
$3$ & $64$ & $1, 6, 15, 20, 15, 6, 1$ \\
$4$ & $256$ & $1, 8, 28, 56, 70, 56, 28, 8, 1$ \\
\bottomrule
\end{tabular}
\end{center}
\end{remark}

% ======================================================================
%
% SUBSECTION: THE VIRASORO AT c = 26
%
% ======================================================================

\subsection{The critical string at $c = 26$}
% label removed: subsec:critical-c26
\index{critical string!c=26@$c = 26$|textbf}
\index{Virasoro!c=26@$c=26$!no-ghost theorem}

We specialize the entire construction to the Virasoro algebra.
The modular characteristic is
\begin{equation}% label removed: eq:g9-kappa-vir
\kappa(\mathrm{Vir}_c) = \frac{c}{2}
\end{equation}
(Theorem~\ref{V1-thm:vir-genus1-curvature}), arising from the quartic
pole $T_{(3)}T = c/2$ in the Virasoro OPE. The Koszul dual is
$\mathrm{Vir}_c^! \simeq \mathrm{Vir}_{26-c}$
(Example~\ref{V1-ex:virasoro-koszul-dual}).

\subsubsection{The Koszul duality pattern}
% label removed: subsubsec:koszul-pattern
\index{Virasoro!Koszul dual}

The complementarity of modular characteristics is
\begin{equation}% label removed: eq:g9-kappa-sum
\kappa(\mathrm{Vir}_c) + \kappa(\mathrm{Vir}_{26-c})
= \frac{c}{2} + \frac{26-c}{2} = 13.
\end{equation}
The secondary anomaly elements for the Koszul pair are
\begin{equation}% label removed: eq:g9-u-pair
u_c = \frac{c}{2} \cdot \omega_g,
\qquad
u_{26-c} = \frac{26-c}{2} \cdot \omega_g.
\end{equation}
Three central charges are structurally distinguished:

\begin{center}
\renewcommand{\arraystretch}{1.4}
\begin{tabular}{clccl}
\toprule
$c$ & Koszul dual & $\kappa$ & $\kappa'$ & Structural property \\
\midrule
$0$ & $\mathrm{Vir}_0^! \simeq \mathrm{Vir}_{26}$ & $0$ & $13$ &
 Uncurved; gravitational regime \\
$13$ & $\mathrm{Vir}_{13}^! \simeq \mathrm{Vir}_{13}$ & $13/2$ & $13/2$ &
 Self-dual; symmetric anomaly \\
$26$ & $\mathrm{Vir}_{26}^! \simeq \mathrm{Vir}_0$ & $13$ & $0$ &
 Maximally curved; dual is gravitational \\
\bottomrule
\end{tabular}
\end{center}

\subsubsection{Two independent vanishing conditions}
% label removed: subsubsec:effective-curvature
\index{complementarity asymmetry|textbf}

The Virasoro family at central charge $c$ is controlled by two
independent scalar conditions. Each governs a different
structural phenomenon; neither implies the other.

\begin{definition}[Complementarity asymmetry]
% label removed: def:effective-curvature
\index{complementarity asymmetry!definition}
For a Koszul pair $(\cA, \cA^!)$ with modular characteristics
$\kappa := \kappa(\cA)$ and $\kappa' := \kappa(\cA^!)$, the
\emph{complementarity asymmetry} is
\begin{equation}% label removed: eq:g9-kappa-eff
\delta_\kappa \;:=\; \kappa - \kappa'.
\end{equation}
Its vanishing is the condition for Koszul self-duality:
$\delta_\kappa = 0$ if and only if $\kappa = \kappa'$.
\end{definition}

For the Virasoro, $\delta_\kappa = c/2 - (26-c)/2 = c - 13$.

The complementarity asymmetry is intrinsic to the Koszul
pair. It measures how unevenly the anomaly budget is
distributed between $\cA$ and $\cA^!$. The second
invariant is extrinsic: it belongs to the physical
matter--ghost system, not to the Koszul pair.

\begin{definition}[Effective curvature]
\index{effective curvature!definition}
The \emph{effective curvature} of a matter--ghost coupling is
\[
\kappa_{\mathrm{eff}} \;:=\;
\kappa(\mathrm{matter}) + \kappa(\mathrm{ghost}),
\]
where the ghost contribution comes from the physical ghost
system (the $bc$ system in the Virasoro case, with
$\kappa_{bc} = -13$ at all central charges). Its vanishing
is anomaly cancellation.
\end{definition}

For the Virasoro at central charge $c$,
$\kappa_{\mathrm{eff}} = c/2 + (-13) = (c-26)/2$. The
effective curvature is \emph{not} the sum
$\kappa + \kappa' = c/2 + (26-c)/2 = 13$, which is a
topological invariant of the Koszul pair (it never vanishes).

The two conditions are related by
$\delta_\kappa = 2\kappa_{\mathrm{eff}} + 13$ and vanish
at different central charges, $13$ units apart:
\[
\delta_\kappa = 0
\;\;\Longleftrightarrow\;\;
c = 13
\quad\text{(Koszul self-duality)},
\qquad
\kappa_{\mathrm{eff}} = 0
\;\;\Longleftrightarrow\;\;
c = 26
\quad\text{(anomaly cancellation)}.
\]

\begin{remark}[Family dependence]
% label removed: rem:g9-complementarity-constraint
The complementarity sum $\kappa + \kappa'$ depends on the
family. For affine Kac--Moody algebras,
$\kappa + \kappa' = 0$
(Remark~\ref{V1-rem:vir-vs-km-complementarity}):
self-duality forces $\kappa = 0$, and if one identifies the
ghost system with the Koszul dual then
$\kappa_{\mathrm{eff}} = \kappa + \kappa' = 0$
identically. For the Virasoro, $\kappa + \kappa' = 13$:
self-duality distributes the anomaly symmetrically at
$\kappa = \kappa' = 13/2$, and anomaly cancellation
requires the separate bc ghost mechanism.
\end{remark}


\begin{remark}[The de Sitter side of Koszul duality]
\label{rem:ads-ds-koszul}
\index{Koszul duality!AdS/dS interpretation}
\index{de Sitter!negative kappa regime}
The complementarity identity
$\kappa(\mathrm{Vir}_c) + \kappa(\mathrm{Vir}_{26-c}) = 13$
admits a gravitational reading. When $c < 26$, both
$\kappa(\mathrm{Vir}_c) = c/2$ and
$\kappa(\mathrm{Vir}_{26-c}) = (26-c)/2$ are positive:
the Koszul pair lives entirely on the anti-de~Sitter side
of the anomaly budget. At $c = 26$ the dual characteristic
$\kappa(\mathrm{Vir}_{26-c}) = (26-c)/2$ vanishes;
the Koszul dual is uncurved, and the matter system
absorbs the full anomaly $\kappa(\mathrm{Vir}_{26}) = 13$.

For $c > 26$ the Koszul dual has
$\kappa(\mathrm{Vir}_{26-c}) = (26-c)/2 < 0$:
the secondary anomaly $u_{26-c} = \kappa(\mathrm{Vir}_{26-c})\cdot\omega_g$
reverses sign, and the fiberwise curvature
$\dfib^{\,2} = \kappa(\mathrm{Vir}_{26-c})\cdot\omega_g$
on the Koszul-dual side becomes negative.
This is the de~Sitter regime of Koszul duality:
the dual algebra $\mathrm{Vir}_{26-c}^{}$
carries a curvature of cosmological sign.

The self-dual point $c = 13$ sits at the transition.
There $\kappa(\mathrm{Vir}_{13}) = \kappa(\mathrm{Vir}_{13}^!) = 13/2$,
and the anomaly budget is distributed symmetrically
between the algebra and its Koszul dual.
The complementarity sum $\kappa + \kappa^{\prime} = 13$
is the total curvature shared by the AdS and dS sectors;
at $c = 13$ the boundary between them passes through
the fixed point of the Feigin--Frenkel involution.
\end{remark}
\subsubsection{The dichotomy theorem}
% label removed: subsubsec:dichotomy-theorem

\begin{theorem}[Critical string dichotomy for the Virasoro;
\ClaimStatusProvedHere]
% label removed: thm:critical-string-dichotomy
\index{critical string!dichotomy theorem|textbf}
The genus-$g$ Clifford completion $G_g(\mathrm{Vir}_c)$ realizes two
regimes:

\begin{enumerate}[label=\textup{(\roman*)}]
\item \emph{Topological regime} ($c \neq 0$):\;
$u_c = (c/2) \cdot \omega_g \neq 0$. The Clifford completion is
Morita equivalent to the transgression algebra:
\[
G_g(\mathrm{Vir}_c)
\simeq_{\mathrm{Morita}}
\mathrm{Mat}_{2^g}\bigl(B_\Theta^{(g)}(\mathrm{Vir}_c)[u_c^{-1}]\bigr).
\]
The genus-$g$ free energy is
\begin{equation}% label removed: eq:g9-partition-topo
F_g(\mathrm{Vir}_c)
= \frac{c}{2} \cdot \lambda_g^{\mathrm{FP}},
\end{equation}
recovering Volume~I Theorem~\ref{V1-thm:vir-all-genera}.

\item \emph{Gravitational regime} ($c = 0$):\;
$u_0 = 0$. The Clifford completion splits:
\[
G_g(\mathrm{Vir}_0)
\cong
\barB^{(g)}(\mathrm{Vir}_0)
\otimes \bigwedge\nolimits^*(\Bbbk^{2g}).
\]
The bar complex is a genuine cochain complex ($\dfib^{\,2} = 0$)
and the exterior factor is the ghost sector.

\item \emph{The $c = 26$ algebra}\/ is in the topological regime
($\kappa = 13 \neq 0$), but its Koszul dual $\mathrm{Vir}_0$ is in
the gravitational regime. The combined Clifford completion is
\begin{equation}% label removed: eq:g9-combined-clifford
G_g(\mathrm{Vir}_{26}) \otimes G_g(\mathrm{Vir}_0)
\;\cong\;
G_g(\mathrm{Vir}_{26})
\otimes \barB^{(g)}(\mathrm{Vir}_0)
\otimes \bigwedge\nolimits^*(\Bbbk^{2g}).
\end{equation}
The full genus-$g$ obstruction budget is carried by
$\mathrm{Vir}_{26}$; the dual $\mathrm{Vir}_0$ contributes only the
exterior algebra factor.
\end{enumerate}
\end{theorem}

\begin{proof}
Parts (i) and (ii) are direct applications of
Volume~I Theorems~\ref{V1-thm:topological-regime}
and~\ref{V1-thm:gravitational-regime} to $\cA = \mathrm{Vir}_c$ with
$\kappa = c/2$. The free energy
formula~\eqref{V1-eq:g9-partition-topo} is Volume~I Theorem~\ref{V1-thm:vir-all-genera}.

For (iii): $\mathrm{Vir}_{26}$ has $\kappa = 13 \neq 0$
(topological); $\mathrm{Vir}_0$ has $\kappa = 0$ (gravitational).
The tensor product follows from K\"unneth. The
complementarity statement is
Volume~I Proposition~\ref{V1-prop:vir-complementarity}: $F_g(\mathrm{Vir}_{26}) +
F_g(\mathrm{Vir}_0) = 13 \cdot \lambda_g^{\mathrm{FP}} + 0 =
13 \cdot \lambda_g^{\mathrm{FP}}$.
\end{proof}

\begin{remark}[Ghost zero modes in bar-cobar language]
% label removed: rem:g9-no-ghost
\index{no-ghost theorem!bar-cobar interpretation}
The classical no-ghost theorem states that the physical states of the
$c = 26$ bosonic string are transverse oscillators. In the present
framework, the Koszul dual $\mathrm{Vir}_{26}^! \simeq \mathrm{Vir}_0$
is in the gravitational regime: $u_0 = 0$, so the handle part of the
dual ghost sector is the exterior algebra
$\bigwedge^*(\Bbbk^{2g})$ with no Clifford entanglement, and
$\barB^{(g)}(\mathrm{Vir}_0)$ is a genuine cochain complex with no
scalar curvature correction. This gives an algebraic model for the
ghost zero-mode count. The passage from this exterior handle sector to
transverse physical states is the Goddard--Thorn and
Frenkel--Garland--Zuckerman spectral-sequence collapse, recorded below
as theorem data in the Clifford-degeneration statement. The ghost
oscillators are the $b$-$c$ modes; the standard zero-mode relation
$\{b_0,c_0\}=1$ supplies the BRST contracting pair, while the
$u_0=0$ Clifford degeneration supplies the exterior handle algebra.

The ghost-number anomaly $2g$ is the count of ghost zero modes
on $\Sigma_g$. The exterior (vs.\ Clifford) structure of the ghost
sector is equivalent to the vanishing of the dual's curvature.
\end{remark}

\subsubsection{Clifford data for the Virasoro at all genera}
% label removed: subsubsec:vir-clifford-data

\begin{computation}[Genus-$g$ Virasoro Clifford data]
% label removed: comp:g9-vir-clifford-data
\index{Virasoro!Clifford completion data}

\begin{center}
\renewcommand{\arraystretch}{1.3}
{\small
\begin{tabular}{cccccl}
\toprule
$c$ & $\kappa$ & $u$ & $\mathrm{Cl}_{2g}(u)$ &
$\dim$ & Regime \\
\midrule
$0$ & $0$ & $0$ &
$\bigwedge^*(\Bbbk^{2g})$ & $2^{2g}$ &
Gravitational \\
$1$ & $1/2$ & $(1/2)\omega_g$ &
$\mathrm{Mat}_{2^g}(\Bbbk)$ & $2^{2g}$ &
Topological \\
$13$ & $13/2$ & $(13/2)\omega_g$ &
$\mathrm{Mat}_{2^g}(\Bbbk)$ & $2^{2g}$ &
Topological (self-dual) \\
$26$ & $13$ & $13\omega_g$ &
$\mathrm{Mat}_{2^g}(\Bbbk)$ & $2^{2g}$ &
Topological (dual: gravitational) \\
\bottomrule
\end{tabular}
}
\end{center}

\medskip\noindent
Genus-$g$ free energy at selected central charges:
\begin{center}
\renewcommand{\arraystretch}{1.3}
{\small
\begin{tabular}{ccccc}
\toprule
$c$ & $F_1$ & $F_2$ & $F_3$ &
$F_g$ \\
\midrule
$0$ & $0$ & $0$ & $0$ & $0$ \\
$1$ & $\frac{1}{48}$ & $\frac{7}{11520}$ &
$\frac{31}{1935360}$ &
$\frac{1}{2}\lambda_g^{\mathrm{FP}}$ \\
$13$ & $\frac{13}{48}$ & $\frac{91}{11520}$ &
$\frac{403}{1935360}$ &
$\frac{13}{2}\lambda_g^{\mathrm{FP}}$ \\
$26$ & $\frac{13}{24}$ & $\frac{91}{5760}$ &
$\frac{403}{967680}$ &
$13\,\lambda_g^{\mathrm{FP}}$ \\
\bottomrule
\end{tabular}
}
\end{center}
where $\lambda_g^{\mathrm{FP}} = \frac{2^{2g-1}-1}{2^{2g-1}} \cdot
\frac{|B_{2g}|}{(2g)!}$. These agree with Volume~I
Theorem~\ref{V1-thm:vir-all-genera} and Volume~I
Computation~\ref{V1-comp:vir-physical-cc}.
\end{computation}

% ======================================================================
%
% SUBSECTION: THE GHOST SECTOR AT c = 26
%
% ======================================================================

\subsection{The exterior algebra structure at $c = 26$}
% label removed: subsec:exterior-c26
\index{critical string!exterior algebra structure}

At $c = 26$, the algebra itself is in the topological regime, but its
\emph{dual} $\mathrm{Vir}_0$ is in the gravitational regime. The ghost
sector of the combined Koszul pair is an exterior algebra.

\begin{theorem}[Ghost sector at $c = 26$; \ClaimStatusProvedHere]
% label removed: thm:ghost-sector-c26
\index{ghost sector!c=26@$c=26$|textbf}
For the Koszul dual $\mathrm{Vir}_{26}^! \simeq \mathrm{Vir}_0$
at genus~$g$:
\begin{enumerate}[label=\textup{(\roman*)}]
\item The Clifford completion of $\mathrm{Vir}_0$ is
\begin{equation}% label removed: eq:g9-ghost-sector
G_g(\mathrm{Vir}_0)
\cong \barB^{(g)}(\mathrm{Vir}_0) \otimes
\bigwedge\nolimits^*\bigl(H_1(\Sigma_g, \Bbbk)\bigr).
\end{equation}

\item The generators are identified with $H_1$-classes:
$\alpha_k \leftrightarrow [A_k]$, $\beta_k \leftrightarrow [B_k]$.

\item The top exterior power
$\bigwedge^{2g}(H_1(\Sigma_g)) \cong \Bbbk$ is generated by
$\mathrm{vol}_g = \alpha_1\beta_1 \cdots \alpha_g\beta_g$.

\item The ghost-number operator
$N_{\mathrm{gh}} = \sum_{k=1}^g(\alpha_k\partial_{\alpha_k} +
\beta_k\partial_{\beta_k})$ has eigenvalues $0, 1, \ldots, 2g$.
The ghost-number anomaly is $2g$: the path-integral measure requires
insertion of $\mathrm{vol}_g$ to produce a nonzero result.
\end{enumerate}
\end{theorem}

\begin{proof}
\emph{(i)} is Theorem~\ref{V1-thm:gravitational-regime}(iii) applied to
$\mathrm{Vir}_0$ (where $\kappa = 0$ and $\lambda = 0$).

\emph{(ii)} The handle operators are fiber integrals of $\eta^{(g)}$
over cycles. When $u = 0$, these satisfy purely exterior algebra
relations, identifying them with $H_1(\Sigma_g, \Bbbk)$.

\emph{(iii)} $\bigwedge^{2g}(H_1) \cong \det(H_1)$ is one-dimensional.

\emph{(iv)} The ghost-number operator is the standard
$\bigwedge$-grading.
\end{proof}

\begin{remark}[Comparison with BRST ghost counting]
% label removed: rem:g9-brst-comparison
\index{BRST cohomology!c=26@$c=26$ comparison}
In the BRST approach, the ghost system has $c_{\mathrm{ghost}} = -26$,
giving total $c_{\mathrm{matter}} + c_{\mathrm{ghost}} = 0$. The
genus-$g$ ghost determinant produces the factor
$\bigwedge^{3g-3}(H^0(\Sigma_g, K^2))$ in the path-integral measure.
The bar-cobar ghost sector is $\bigwedge^{2g}(H_1(\Sigma_g))$, not
$\bigwedge^{3g-3}$. The discrepancy arises from different
gradings: bar-cobar ghosts count \emph{homological} cycles
(intersection form on $H_1$); BRST ghosts count \emph{conformal}
zero modes ($b$-ghost on $H^0(K^2)$). The Riemann--Roch relation
$\dim H^0(K^2) = 3g-3$ (for $g \geq 2$) connects the two counts
via Serre duality and the period matrix.
\end{remark}

% ======================================================================
%
% SUBSECTION: SELF-DUALITY AT c = 13
%
% ======================================================================

\subsection{Self-duality at $c = 13$}
% label removed: subsec:self-duality-c13
\index{Virasoro!c=13@$c=13$!self-duality|textbf}

The Feigin--Frenkel involution $c \mapsto 26 - c$ has its fixed point
at $c = 13$. At this central charge, the Virasoro is Koszul self-dual
in the \emph{curved} sense.

\begin{theorem}[Virasoro self-duality at $c = 13$;
\ClaimStatusProvedHere]
% label removed: thm:vir-self-duality-c13
\index{Virasoro!c=13@$c=13$!self-duality theorem}
At $c = 13$:
\begin{enumerate}[label=\textup{(\roman*)}]
\item $\mathrm{Vir}_{13}^! \simeq \mathrm{Vir}_{13}$: Koszul
self-duality.

\item $\kappa(\mathrm{Vir}_{13}) = 13/2$: symmetric anomaly budget,
$\kappa = \kappa' = 13/2$.

\item The secondary anomalies coincide:
$u_{13} = u'_{13} = (13/2) \cdot \omega_g$.

\item The Clifford completions are isomorphic:
$G_g(\mathrm{Vir}_{13}) \cong G_g(\mathrm{Vir}_{13}^!)$, both
topological.

\item Symmetric complementarity:
\begin{equation}% label removed: eq:g9-c13-complementarity
F_g(\mathrm{Vir}_{13}) = F_g(\mathrm{Vir}_{13}^!)
= \frac{13}{2} \cdot \lambda_g^{\mathrm{FP}},
\end{equation}
with total $F_g + F_g' = 13 \cdot \lambda_g^{\mathrm{FP}}$.

\item Complementarity asymmetry vanishes: $\delta_\kappa = 0$.
\end{enumerate}
\end{theorem}

\begin{proof}
\emph{(i)} $\mathrm{Vir}_c^! \simeq \mathrm{Vir}_{26-c}$ at $c = 13$
gives $\mathrm{Vir}_{13}^! \simeq \mathrm{Vir}_{13}$.
\emph{(ii)--(vi)} follow from (i) and the definitions.
\end{proof}

\begin{remark}[Two self-duality points]
% label removed: rem:g9-two-self-dualities
\index{Virasoro!two self-dualities}
The Virasoro has two distinct self-duality points:
\begin{enumerate}[label=\textup{(\alph*)}]
\item \emph{Quadratic self-duality at $c = 0$}
(Theorem~\ref{V1-thm:virasoro-self-duality}). Here $\kappa = 0$: the bar
complex is a genuine cochain complex, self-duality is of \emph{classical}
(quadratic) type ($\mathrm{Vir}_0^{!_{\mathrm{quad}}} \simeq \mathrm{Vir}_0$),
and the Clifford completion is in the gravitational
regime. The chiral Koszul dual remains
$\mathrm{Vir}_0^! = \mathrm{Vir}_{26} \neq \mathrm{Vir}_0$.

\item \emph{Chiral Koszul self-duality at $c = 13$}
(Theorem~\ref{V1-thm:vir-self-duality-c13}). Here $\kappa = 13/2 \neq 0$:
the bar complex is curved, self-duality is of \emph{modular}
(Feigin--Frenkel) type, and the Clifford completion is in the topological
regime.
\end{enumerate}
At $c = 0$, ghosts are free (exterior algebra); at $c = 13$, ghosts
are entangled (Clifford algebra with $u = (13/2)\omega_g \neq 0$).
The two self-dualities are genuinely different: $c = 0$ kills the
curvature; $c = 13$ distributes it symmetrically.
\end{remark}

% ======================================================================
%
% SUBSECTION: THE DICHOTOMY SUMMARY
%
% ======================================================================

\subsection{The critical dichotomy: structural summary}
% label removed: subsec:dichotomy-summary
\index{critical string!dichotomy!summary}

\begin{theorem}[Critical string dichotomy: structural summary;
\ClaimStatusProvedHere]
% label removed: thm:critical-dichotomy-summary
\index{critical string!dichotomy!structural summary}
The genus-$g$ algebraic structure of the Virasoro Koszul pair
$(\mathrm{Vir}_c, \mathrm{Vir}_{26-c})$ is controlled by
$\kappa = c/2$, $\kappa' = (26-c)/2$,
$\delta_\kappa = c - 13$, and falls into the following
pattern:

\begin{center}
\renewcommand{\arraystretch}{1.4}
\begin{tabular}{lccc}
\toprule
& $c < 13$ & $c = 13$ & $c > 13$ \\
\midrule
$\delta_\kappa$ & $< 0$ & $= 0$ & $> 0$ \\
Dominant anomaly & Dual $\mathrm{Vir}_{26-c}$ &
Equal & $\mathrm{Vir}_c$ \\
Self-dual? & No & Yes & No \\
Both regimes & Topological & Topological & Topological \\
\bottomrule
\end{tabular}
\end{center}

\smallskip\noindent
The two boundary cases:
\begin{center}
\renewcommand{\arraystretch}{1.4}
\begin{tabular}{lcc}
\toprule
& $c = 0$ & $c = 26$ \\
\midrule
$\delta_\kappa$ & $-13$ & $+13$ \\
$\mathrm{Vir}_c$ regime & Gravitational & Topological \\
$\mathrm{Vir}_{26-c}$ regime & Topological & Gravitational \\
Ghost sector & $\bigwedge^*(H_1)$ at $c=0$ &
$\bigwedge^*(H_1)$ at $c'=0$ \\
Anomaly budget & Entirely on dual & Entirely on algebra \\
\bottomrule
\end{tabular}
\end{center}

\smallskip\noindent
The dichotomy is:

\emph{At $c = 26$, the Koszul dual $\mathrm{Vir}_0$ is in the
gravitational regime \textup{(}$u = 0$, exterior algebra ghosts,
uncurved bar complex\textup{)}, while $\mathrm{Vir}_{26}$ itself
carries the entire anomaly budget. At $c = 0$, the roles are
reversed. At $c = 13$, the anomaly is shared equally and the system
is self-dual.}
\end{theorem}

\begin{proof}
Compilation of Volume~I
Theorems~\ref{V1-thm:critical-string-dichotomy},
\ref{V1-thm:ghost-sector-c26},
and~\ref{V1-thm:vir-self-duality-c13}, together with Volume~I
Proposition~\ref{V1-prop:vir-complementarity}.
\end{proof}

% ======================================================================
%
% SUBSECTION: SHADOW TOWER AT CRITICAL POINTS
%
% ======================================================================

\subsection{The shadow obstruction tower at the critical points}
% label removed: subsec:shadow-critical
\index{shadow tower!critical points}

The critical string dichotomy interacts with the shadow obstruction tower
$\Theta_\cA^{\leq r}$ through the modular characteristic
$\kappa = \Theta_\cA^{\leq 2}$.

\begin{proposition}[Shadow obstruction tower at critical central charges;
\ClaimStatusProvedHere]
% label removed: prop:shadow-critical
\index{shadow tower!critical central charges}
\begin{enumerate}[label=\textup{(\roman*)}]
\item \emph{At $c = 0$}: $\kappa = 0$. The quartic pole vanishes, so
the OPE is purely quadratic ($T_{(3)}T = 0$) and the bar complex is
uncurved. The universal genus-$1$ scalar term vanishes, but
$\kappa = 0$ does not by itself determine the higher-degree shadow
tower.

\item \emph{At $c = 13$}: $\kappa = 13/2 \neq 0$. The shadow obstruction tower
is nonzero at every level: $\Theta^{\leq 2} = 13/2$,
$\Theta^{\leq 4}$ involves the quartic resonance class
$Q^{\mathrm{contact}} = 10/[13(5\cdot 13 + 22)] = 10/1131$.
The tower is infinite ($r_{\max} = \infty$, mixed type) and
self-dual: $\Theta_{\mathrm{Vir}_{13}}^{\leq r} =
\Theta_{\mathrm{Vir}_{13}^!}^{\leq r}$ at every level.

\item \emph{At $c = 26$}: $\kappa = 13 \neq 0$. The quartic
contact invariant is
$Q^{\mathrm{contact}} = 10/[26 \cdot 152] = 5/1976$.
The dual algebra $\mathrm{Vir}_0$ lies on the uncurved scalar locus,
so $\Theta_{\mathrm{Vir}_0}^{\min} = 0$ and
$F_1(\mathrm{Vir}_0) = 0$. This scalar datum does not determine the
dual higher-degree shadow obstruction tower. Thus the nonzero scalar and quartic
data exhibited here come from $\mathrm{Vir}_{26}$, while the dual
higher-degree tower is not fixed by $\kappa(\mathrm{Vir}_0)=0$ alone.
\end{enumerate}
\end{proposition}

\begin{proof}
\emph{(i)} $\kappa(\mathrm{Vir}_0) = 0$; at $c = 0$ the quartic pole
vanishes and the OPE is quadratic. Hence the scalar class
$\Theta_{\mathrm{Vir}_0}^{\min}$ and the universal genus-$1$ scalar
term vanish, but no conclusion about the full higher-degree tower
follows from $\kappa(\mathrm{Vir}_0)=0$ alone.

\emph{(ii)} Self-duality implies $\Theta_{\mathrm{Vir}_{13}}^{\leq r}
= \Theta_{\mathrm{Vir}_{13}^!}^{\leq r}$. The quartic invariant
evaluates via $Q^{\mathrm{contact}} = 10/[c(5c+22)]|_{c=13}$.

\emph{(iii)} $\kappa(\mathrm{Vir}_0) = 0$ gives the vanishing of the
dual scalar class $\Theta_{\mathrm{Vir}_0}^{\min}$ and of
$F_1(\mathrm{Vir}_0)$; the higher-degree dual tower is not determined
by this scalar input alone.
\end{proof}

\begin{remark}[Depth-zero resonance shadow revisited]
% label removed: rem:g9-depth-zero-revisited
\index{depth-zero resonance shadow!revisited}
The same-family shadow partner $\mathrm{Vir}_{26-c}$ controls
M/S-level complementarity
(Remark~\ref{V1-rem:virasoro-resonance-model}). At $c = 26$, this
partner is $\mathrm{Vir}_0$, the dual algebra on the uncurved scalar
locus, not a trivial algebra. The depth-zero resonance shadow
saturates only at the scalar/genus-$1$ level; higher-degree data of the
dual partner are not fixed by that cancellation. This is the
algebraic content of the critical string: the universal scalar budget
closes at $c = 26$, while the non-scalar shadow obstruction tower still requires
separate control.
\end{remark}

% ======================================================================
%
% SUBSECTION: HODGE-CLIFFORD COMPATIBILITY
%
% ======================================================================

\subsection{The Hodge-theoretic content}
% label removed: subsec:hodge-dichotomy
\index{Hodge theory!critical string dichotomy}

\begin{proposition}[Hodge--Clifford compatibility;
\ClaimStatusProvedHere]
% label removed: prop:hodge-clifford
\index{Hodge theory!Clifford compatibility}
The Clifford algebra $\mathrm{Cl}_{2g}(u)$ is compatible with the
Hodge decomposition $H^1(\Sigma_g, \C) = H^{1,0} \oplus H^{0,1}$:
\begin{enumerate}[label=\textup{(\roman*)}]
\item The $A$-cycle operators $\alpha_k$ correspond to
$H^{1,0}(\Sigma_g)$ (holomorphic differentials).

\item The $B$-cycle operators $\beta_k$ correspond to
$H^{0,1}(\Sigma_g)$ (antiholomorphic differentials).

\item The Clifford relation $\alpha_i \beta_j + \beta_j \alpha_i =
\delta_{ij}\, u$ encodes the Hodge--Riemann relation:
\[
\int_{\Sigma_g} \omega_i \wedge \bar\omega_j = \delta_{ij}\, \omega_g,
\]
where $\{\omega_1, \ldots, \omega_g\}$ are the normalized abelian
differentials ($\oint_{A_l}\omega_k = \delta_{kl}$).

\item The period matrix $\Omega_{kl} = \oint_{B_l}\omega_k$ appears as
the overlap $\langle \alpha_k, \beta_l \rangle = \Omega_{kl}$ in the
complexified Clifford algebra.
\end{enumerate}
\end{proposition}

\begin{proof}
\emph{(i)--(ii)} The symplectic basis $(A_k, B_k)$ is Poincar\'e dual
to the Hodge basis $(\omega_k, \bar\omega_k)$.

\emph{(iii)} The Riemann bilinear relation gives
$\frac{i}{2}\int_{\Sigma_g}\omega_k \wedge \bar\omega_l =
(\operatorname{Im}\Omega)_{kl}$ in the Arakelov normalization.
Integration against $u = \kappa \cdot \omega_g$ yields
$\kappa \cdot \delta_{kl} = \delta_{kl}\, u$, matching the
Clifford relation.

\emph{(iv)} $\Omega_{kl} = \oint_{B_l}\omega_k$ determines the
$\alpha_k$--$\beta_l$ overlap in the complexification.
\end{proof}

\begin{remark}[The Arakelov metric as Clifford trace]
% label removed: rem:g9-arakelov-clifford
\index{Arakelov metric!Clifford trace}
The Arakelov form $\omega_g = \frac{i}{2}\sum_{k=1}^g
\omega_k \wedge \bar\omega_k$
(Proposition~\ref{V1-prop:genus-g-curvature-package}) is the trace of the
Clifford form: $\omega_g = \frac{1}{2g}\sum_{k=1}^g
(\alpha_k\beta_k + \beta_k\alpha_k)|_{u=1}$.
\end{remark}

% ======================================================================
%
% SUBSECTION: CLIFFORD SPECTRAL SEQUENCE
%
% ======================================================================

\subsection{The spectral sequence of the Clifford filtration}
% label removed: subsec:clifford-spectral
\index{spectral sequence!Clifford filtration}

\begin{construction}[Clifford spectral sequence]
% label removed: constr:clifford-spectral
\index{spectral sequence!Clifford|textbf}
The Clifford degree filtration on $G_g(\cA)$,
\begin{equation}% label removed: eq:g9-clifford-filtration
F^p G_g := B_\Theta^{(g)}(\cA) \otimes
\bigoplus_{k \leq p} \mathrm{Cl}_{2g}^{(k)}(u),
\end{equation}
where $\mathrm{Cl}_{2g}^{(k)}$ is the span of products of at most $k$
generators, is compatible with $d_{G_g}$ (since handle operators are
closed). The $E_1$ page is
\begin{equation}% label removed: eq:g9-clifford-E1
E_1^{p,q} = H^{p+q}(B_\Theta^{(g)}(\cA))
\otimes \mathrm{gr}^p_{\mathrm{Cl}}\,\mathrm{Cl}_{2g}(u).
\end{equation}
\end{construction}

\begin{proposition}[Degeneration; \ClaimStatusProvedHere]
% label removed: prop:clifford-spectral-degeneration
\begin{enumerate}[label=\textup{(\roman*)}]
\item \emph{Gravitational regime} ($u = 0$): degeneration at $E_1$;
$\mathrm{gr}_{\mathrm{Cl}}^* = \bigwedge^*(\Bbbk^{2g})$; hence
$H^*(G_g|_{u=0}) \cong H^*(B) \otimes \bigwedge^*(\Bbbk^{2g})$.

\item \emph{Topological regime} ($u$ invertible): degeneration at
$E_1$; $\mathrm{gr}_{\mathrm{Cl}}^*$ is the associated graded of a
matrix algebra; $H^*(G_g[u^{-1}]) \cong
\mathrm{Mat}_{2^g}(H^*(B_\Theta[u^{-1}]))$.
\end{enumerate}
\end{proposition}

\begin{proof}
In both cases, the handle operators are closed and the Clifford
multiplication preserves the filtration, so $d_r = 0$ for $r \geq 1$.
\end{proof}

% ======================================================================
%
% SUBSECTION: BOSONIC STRING IDENTIFICATIONS
%
% ======================================================================

\subsection{The $c = 26$ algebra and the bosonic string}
% label removed: subsec:bosonic-string
\index{bosonic string!bar-cobar interpretation|textbf}

\begin{proposition}[Bosonic string structural identifications;
\ClaimStatusProvedHere]
% label removed: prop:bosonic-string-identifications
\index{bosonic string!structural identifications}
At $c = 26$:
\begin{enumerate}[label=\textup{(\roman*)}]
\item \emph{Anomaly budget.}\;
$F_g(\mathrm{Vir}_{26}) = 13 \cdot \lambda_g^{\mathrm{FP}}$,
$F_g(\mathrm{Vir}_0) = 0$.

\item \emph{Koszul pair.}\;
$(\mathrm{Vir}_{26}, \mathrm{Vir}_0)$ with $\kappa + \kappa' = 13$,
maximally asymmetric: $|\delta_\kappa| = 13$.

\item \emph{Ghost sector.}\;
$\bigwedge^*(H_1(\Sigma_g)) \cong \bigwedge^{2g}$, ghost-number
anomaly $2g$.

\item \emph{Generating function.}\;
$\sum_{g=1}^\infty F_g(\mathrm{Vir}_{26})\,x^{2g}
= 13(x/(2\sin(x/2)) - 1)$
with radius of convergence $|x| = 2\pi$.

\item \emph{Transgression.}\;
$B_\Theta^{(g)}(\mathrm{Vir}_{26})$ has $u = 13\omega_g$
(strongly topological); $B_\Theta^{(g)}(\mathrm{Vir}_0)$ has
$\eta^2 = 0$ (uncurved).

\item \emph{Combined Clifford.}\;
\begin{equation}% label removed: eq:g9-combined-final
G_g(\mathrm{Vir}_{26}) \otimes G_g(\mathrm{Vir}_0)
\cong
\mathrm{Mat}_{2^g}(B_\Theta^{(g)}(\mathrm{Vir}_{26})[u^{-1}])
\otimes \barB^{(g)}(\mathrm{Vir}_0)
\otimes \bigwedge\nolimits^*(\Bbbk^{2g}).
\end{equation}

\item \emph{Physical content.}\; Morita reduction gives
\begin{equation}% label removed: eq:g9-physical-content
D(G_g(\mathrm{Vir}_{26}) \otimes G_g(\mathrm{Vir}_0)
\text{-}\mathsf{mod})
\simeq
D\bigl(B_\Theta[u^{-1}] \otimes \barB(\mathrm{Vir}_0) \otimes
\bigwedge\nolimits^*(\Bbbk^{2g})\text{-}\mathsf{mod}\bigr).
\end{equation}
\end{enumerate}
\end{proposition}

\begin{proof}
All items follow from the established theorems:
(i) Theorem~\ref{V1-thm:vir-all-genera};
(ii) complementarity;
(iii) Theorem~\ref{V1-thm:ghost-sector-c26};
(iv) Remark~\ref{V1-rem:vir-trichotomy-genera};
(v) Definitions~\ref{V1-def:transgression-algebra}
and~\ref{V1-def:bar-transgression-complex};
(vi) Theorems~\ref{V1-thm:topological-regime}
and~\ref{V1-thm:gravitational-regime};
(vii) Theorem~\ref{V1-thm:topological-regime}(iii).
\end{proof}

% ======================================================================
%
% SUBSECTION: KM CONTRAST AND W_N EXTENSION
%
% ======================================================================

\subsection{The affine Kac--Moody contrast}
% label removed: subsec:km-contrast
\index{affine Kac--Moody!critical string contrast}

\begin{remark}[Complementarity sum comparison]
% label removed: rem:g9-km-contrast
\index{complementarity!sum comparison}
\begin{center}
\renewcommand{\arraystretch}{1.3}
\begin{tabular}{lccc}
\toprule
Family & $\kappa + \kappa'$ & Self-dual point &
Gravitational locus \\
\midrule
$\widehat{\fg}_k$ &
$0$ & $k = -h^\vee$ (critical; excluded) &
$\kappa = 0$ iff $k = -h^\vee$ \\
$\mathrm{Vir}_c$ &
$13$ & $c = 13$ &
$\kappa = 0$ iff $c = 0$ \\
$\mathcal{W}_N$ &
$\ne 0$ (generically) & exists for each $N$ &
$\kappa = 0$ at isolated $c$ \\
\bottomrule
\end{tabular}
\end{center}
For affine Kac--Moody at non-critical level, $\kappa + \kappa' = 0$
and $\kappa \neq 0$: there is no gravitational regime. The
critical string dichotomy is intrinsically a Virasoro/$\cW_N$
phenomenon.
\end{remark}

\subsection{The $\mathcal{W}_N$ generalization}
% label removed: subsec:wn-dichotomy
\index{W-algebra@$\mathcal{W}$-algebra!critical dichotomy}

\begin{remark}[$\mathcal{W}_N$ complementarity sums]
% label removed: rem:g9-wn-sums
\index{W-algebra@$\mathcal{W}$-algebra!complementarity sums}
The modular characteristic is $\kappa(\mathcal{W}_{N,c}) = c \cdot
\varrho(\mathfrak{sl}_N)$ where $\varrho = \sum_{i=1}^{N-1}(m_i+1)^{-1}$
is the anomaly ratio.
\begin{center}
\renewcommand{\arraystretch}{1.3}
\begin{tabular}{ccccc}
\toprule
$N$ & $\varrho(\mathfrak{sl}_N)$ & $\kappa$ &
$\kappa + \kappa'$ & Self-dual $c$ \\
\midrule
$2$ & $1/2$ & $c/2$ & $13$ & $13$ \\
$3$ & $5/6$ & $5c/6$ & $250/3$ & $50$ \\
$N$ & $H_N - 1$ & $c(H_N - 1)$ &
$c_{\mathrm{crit}}(H_N-1)$ & $c_{\mathrm{crit}}/2$ \\
\bottomrule
\end{tabular}
\end{center}
Here $H_N = \sum_{k=1}^N 1/k$ is the harmonic number. The
transgression algebra, Clifford completion, and two-regime dichotomy
carry through with $\kappa = c \cdot \varrho$ in place of $c/2$.
\end{remark}

% ======================================================================
%
% SUBSECTION: HOLOGRAPHIC INTERPRETATION
%
% ======================================================================

\subsection{Relation to the holographic datum}
% label removed: subsec:holographic-relation
\index{holographic modular Koszul datum!critical string}

\begin{remark}[Holographic interpretation]
% label removed: rem:g9-holographic
\index{holographic modular Koszul datum!critical string interpretation}
In the holographic setting
(\S\ref{V1-sec:frontier-modular-holography-platonic}):
\begin{enumerate}[label=\textup{(\alph*)}]
\item The \emph{topological regime} ($u$ invertible) corresponds to
the interacting bulk: nontrivial curvature, nonzero shadow obstruction tower,
bulk gravity determined by the Yangian $r(z)$ and the shadow
connection $\nabla^{\mathrm{hol}}$.

\item The \emph{gravitational regime} ($u = 0$) corresponds to the
degenerate bulk: uncurved boundary, vanishing shadow obstruction tower, the bulk
reduces to a topological theory with ghost sector
$\bigwedge^*(H_1(\Sigma_g))$.

\item The \emph{self-dual point} ($\delta_\kappa = 0$)
corresponds to self-conjugacy under gravitational $S$-duality (G4):
the bulk and defect sectors contribute equally.

\item At $c = 26$, the holographic datum
$\mathcal{H}(\mathrm{Vir}_{26})$ is maximally asymmetric: the defect
sector ($\mathrm{Vir}_0$, gravitational) is trivial, and the bulk
sector ($\mathrm{Vir}_{26}$, topological) carries all structure. This
is the bosonic string from the holographic perspective.
\end{enumerate}
\end{remark}

% ======================================================================
%
% SUBSECTION: THE LOCALIZATION SEQUENCE
%
% ======================================================================

\subsection{The localization sequence}
% label removed: subsec:localization-sequence
\index{localization sequence!Clifford completion}

The two regimes are connected by a localization sequence in the
secondary anomaly $u$. This is the algebraic analog of the passage
from generic central charge to the critical point.

\begin{theorem}[Localization sequence; \ClaimStatusProvedHere]
% label removed: thm:localization-sequence
\index{localization sequence!theorem|textbf}
There is a localization fiber sequence of dg algebras:
\begin{equation}% label removed: eq:g9-localization
G_g(\cA)/(u)
\;\longrightarrow\;
G_g(\cA)
\;\longrightarrow\;
G_g(\cA)[u^{-1}],
\end{equation}
with the following properties:
\begin{enumerate}[label=\textup{(\roman*)}]
\item The left term $G_g(\cA)/(u) = \barB^{(g)}(\cA) \otimes
\bigwedge^*(\Bbbk^{2g})$ is the gravitational regime:
the Clifford algebra with $u = 0$.

\item The right term $G_g(\cA)[u^{-1}] \simeq_{\mathrm{Morita}}
\mathrm{Mat}_{2^g}(B_\Theta[u^{-1}])$ is the topological regime.

\item The long exact sequence in cohomology reads:
\begin{equation}% label removed: eq:g9-localization-les
\cdots \to
H^n(G_g/(u)) \xrightarrow{\iota}
H^n(G_g) \xrightarrow{j}
H^n(G_g[u^{-1}]) \xrightarrow{\partial}
H^{n+1}(G_g/(u)) \to \cdots
\end{equation}
The connecting homomorphism $\partial$ measures the failure of
cohomology classes in the topological regime to extend to the
gravitational regime.
\end{enumerate}
\end{theorem}

\begin{proof}
The sequence~\eqref{V1-eq:g9-localization} is the standard localization
for a central element $u$ in a dg algebra: the fiber of
$G_g \to G_g[u^{-1}]$ is $G_g/(u^\infty) = \mathrm{colim}_n\,
G_g/(u^n)$. Since $u$ is central and $u^2 = \kappa^2 \cdot \omega_g^2
\in B$ (where $\omega_g^2$ is a $(2,2)$-class on $\Sigma_g$), the
$u$-torsion stabilizes after finitely many steps. For the Clifford
algebra, setting $u = 0$ in the
relations~\eqref{V1-eq:g9-clifford-ab}--\eqref{V1-eq:g9-clifford-bb} gives
the exterior algebra, confirming (i). Part (ii) is
Theorem~\ref{V1-thm:topological-regime}. Part (iii) is the long exact
sequence of the fiber-cofiber sequence in the dg category.
\end{proof}

\begin{corollary}[Localization for the Virasoro; \ClaimStatusProvedHere]
% label removed: cor:g9-localization-vir
\index{localization sequence!Virasoro}
For the Virasoro, the localization sequence specializes to:
\begin{equation}% label removed: eq:g9-localization-vir
\barB^{(g)}(\mathrm{Vir}_c) \otimes \bigwedge\nolimits^{2g}
\;\longrightarrow\;
G_g(\mathrm{Vir}_c)
\;\longrightarrow\;
\mathrm{Mat}_{2^g}(B_\Theta^{(g)}(\mathrm{Vir}_c)[(c/2)^{-1}]).
\end{equation}
The connecting homomorphism is
\begin{equation}% label removed: eq:g9-connecting-vir
\partial\colon
H^n(\mathrm{Mat}_{2^g}(B_\Theta[(c/2)^{-1}]))
\longrightarrow
H^{n+1}(\barB^{(g)}(\mathrm{Vir}_c))
\otimes \bigwedge\nolimits^{2g}.
\end{equation}
At $c = 0$ the right term vanishes (the localization
$B_\Theta[u^{-1}]$ is trivial when $u = 0$), so the left term
$\barB^{(g)}(\mathrm{Vir}_0) \otimes \bigwedge^{2g}$ accounts for all
cohomology. At $c = 26$ the left term has trivial
differential (uncurved), so the connecting homomorphism
$\partial$ captures the entire obstruction to passing from the
topological regime to the gravitational regime.
\end{corollary}

\begin{proof}
Set $u = (c/2)\omega_g$. The left term is $G_g/(u)$; the right is
$G_g[u^{-1}]$. The statements about $c = 0$ and $c = 26$ follow from
the vanishing of $u$ and the maximality of $\kappa$ respectively.
\end{proof}

% ======================================================================
%
% SUBSECTION: PRIMITIVITY AND BIALGEBRA STRUCTURE
%
% ======================================================================

\subsection{Primitivity of the secondary anomaly}
% label removed: subsec:primitivity
\index{secondary anomaly!primitivity}
\index{bialgebra!transgression algebra}

The bar coalgebra $\barB^{(g)}(\cA)$ carries a coalgebra structure
from the deconcatenation coproduct. The transgression element $\eta$
and the secondary anomaly $u = \eta^2$ interact with this coproduct in
a controlled way.

\begin{proposition}[Bialgebra extension; \ClaimStatusProvedHere]
% label removed: prop:g9-bialgebra
\index{bialgebra!transgression extension}
The transgression algebra $B_\Theta^{(g)}(\cA)$ carries a unique
dg bialgebra structure extending the coalgebra on
$\barB^{(g)}(\cA)$, determined by
\begin{equation}% label removed: eq:g9-eta-coproduct
\Delta(\eta) = \eta \otimes 1 + 1 \otimes \eta,
\qquad
\varepsilon(\eta) = 0.
\end{equation}
In particular, $\eta$ is primitive.
\end{proposition}

\begin{proof}
The coproduct~\eqref{V1-eq:g9-eta-coproduct} is forced by the
requirement that $\Delta$ be an algebra map: $\Delta(\eta^2) =
\Delta(\eta)^2 = (\eta \otimes 1 + 1 \otimes \eta)^2 =
\eta^2 \otimes 1 + \eta \otimes \eta + \eta \otimes \eta +
1 \otimes \eta^2 = \lambda \otimes 1 + 2\eta \otimes \eta +
1 \otimes \lambda$. Here $\lambda = \kappa \cdot \omega_g$ is
the curvature class on $\overline{\cM}_g$, not a fibrewise
coalgebra element: it lives in $H^{2}(\overline{\cM}_g)$ via
the Hodge bundle, and its primitivity in $B_\Theta^{(g)}(\cA)$
is the statement that the
period-corrected differential
$\Dg{g}$ of Volume~I,
Remark~\ref{V1-rem:sc-higher-genus}\textup{(iii)}, is a
coderivation of~$\Delta$. The fibrewise differential
$\dfib$ is \emph{not} a coderivation
(Volume~I,
Remark~\ref{V1-rem:sc-higher-genus}\textup{(ii)}); only after
the Fay trisecant period correction does the curved
$E_1$ dg coalgebra structure persist, and only then does
the Hodge curvature class push forward to a primitive element
of $B_\Theta^{(g)}(\cA)$. The counit $\varepsilon(\eta) = 0$
follows from $\varepsilon(\lambda) = 0$. Coassociativity:
$(\Delta \otimes \id)\Delta(\eta) = \eta \otimes 1 \otimes 1 +
1 \otimes \eta \otimes 1 + 1 \otimes 1 \otimes \eta =
(\id \otimes \Delta)\Delta(\eta)$.
\end{proof}

\begin{corollary}[Primitivity of the secondary anomaly;
\ClaimStatusProvedHere]
% label removed: cor:g9-u-coproduct
The secondary anomaly element is primitive:
\begin{equation}% label removed: eq:g9-u-coproduct
\Delta(u) = u \otimes 1 + 1 \otimes u.
\end{equation}
\end{corollary}

\begin{proof}
Direct computation from Proposition~\ref{V1-prop:g9-bialgebra}.
Using $\Delta(\eta) = \eta \otimes 1 + 1 \otimes \eta$ and the
graded multiplication in $B_\Theta \otimes B_\Theta$
(Koszul sign rule: $(a \otimes b)(c \otimes d)
= (-1)^{|b||c|}\,ac \otimes bd$):
\begin{align*}
\Delta(u) &= \Delta(\eta)^2
= (\eta \otimes 1 + 1 \otimes \eta)^2 \\
&= \eta^2 \otimes 1
 + (\eta \otimes 1)(1 \otimes \eta)
 + (1 \otimes \eta)(\eta \otimes 1)
 + 1 \otimes \eta^2 \\
&= u \otimes 1
 + \eta \otimes \eta
 + (-1)^{|\eta||\eta|}\eta \otimes \eta
 + 1 \otimes u \\
&= u \otimes 1
 + \eta \otimes \eta
 - \eta \otimes \eta
 + 1 \otimes u \\
&= u \otimes 1 + 1 \otimes u.
\end{align*}
The cross terms cancel because $|\eta| = 1$ is odd:
$(1 \otimes \eta)(\eta \otimes 1) = (-1)^{|\eta||\eta|}\eta \otimes \eta
= -\eta \otimes \eta$.
\end{proof}

\begin{remark}[Primitivity and the Milnor--Moore theorem]
% label removed: rem:g9-milnor-moore
\index{Milnor--Moore theorem!transgression}
When $B_\Theta$ is connected and locally finite, the Milnor--Moore
theorem identifies the primitive elements of $H^*(B_\Theta)$ with the
indecomposables of the associated Lie algebra. Since $u$ is primitive
(Corollary~\ref{V1-cor:g9-u-coproduct}), its cohomology class
$[u] \in \mathrm{Prim}(H^*(B_\Theta))$ generates a one-dimensional
abelian sub-Lie algebra. Its vanishing ($\kappa = 0$) is the
condition for the Lie algebra to lose this generator, the algebraic
avatar of the critical string.
\end{remark}

% ======================================================================
%
% SUBSECTION: THE TRANSGRESSION AS UNIVERSAL PROPERTY
%
% ======================================================================

\subsection{The universal property of the transgression algebra}
% label removed: subsec:universal-property
\index{transgression algebra!universal property}

\begin{proposition}[Universal property; \ClaimStatusProvedHere]
% label removed: prop:g9-universal
\index{transgression algebra!universal property|textbf}
Let $(B, d, \lambda)$ be a curved dg algebra with
$d^2 = [\lambda, -]$. For any dg algebra $(C, d_C)$ with
$d_C^2 = 0$ and any dg algebra map $f\colon B \to C$ such that
$f(\lambda) = \xi^2$ for some $\xi \in C^1$ with $d_C(\xi) = 0$,
there exists a unique dg algebra map $\tilde{f}\colon B_\Theta \to C$
extending $f$ with $\tilde{f}(\eta) = \xi$:
\[
\begin{tikzcd}
B \arrow[r, hook] \arrow[d, "f"'] &
B_\Theta \arrow[dl, dashed, "\exists!\,\tilde{f}"] \\
C &
\end{tikzcd}
\]
\end{proposition}

\begin{proof}
Define $\tilde{f}(b_0 + b_1\eta) = f(b_0) + f(b_1)\xi$ for
$b_0, b_1 \in B$. This is well-defined on $B_\Theta$ because
$\tilde{f}(\eta^2 - \lambda) = \xi^2 - f(\lambda) = 0$.
It is an algebra map: $\tilde{f}$ preserves the
multiplication~\eqref{V1-eq:g9-BT-multiplication} because $f$ is an
algebra map and $\xi$ centralizes $\im(f)$ (since $f(\lambda) = \xi^2$
is central in $C$, and $\xi$ commutes with $\im(f)$ by the
same centrality argument used for $\eta$). It respects the
differential: $\tilde{f}(d_{B_\Theta}(b)) = \tilde{f}(d(b) +
[\eta, b]) = f(d(b)) + [f(\eta), f(b)] = d_C(f(b)) +
[\xi, f(b)] = d_C(\tilde{f}(b))$ using $d_C \circ f = f \circ d$
(since $f$ is a dg map, with the correction that $d^2 = [\lambda, -]$
maps to $d_C^2 = 0$ via $f(\lambda) = \xi^2$). Uniqueness: any
extension must send $\eta \mapsto \xi$ (there is no choice), and
the value on $B$ is fixed by $f$.
\end{proof}

\begin{corollary}[The comparison map is universal;
\ClaimStatusProvedHere]
% label removed: cor:g9-comparison-universal
The comparison map $B_\Theta^{(g)}(\cA) \to (\barB^{(g)}(\cA),
\Dg{g})$ of Remark~\ref{V1-rem:g9-mc-relation} is the unique factorization
of the identity on $\barB^{(g)}(\cA)$ through the transgression
algebra, with $\eta$ sent to the period-dependent element
$\sum_k t_k \omega_k$ whose square is $\kappa \cdot \omega_g$ (by the
Lagrangian property of the $A$-cycle subspace).
\end{corollary}

% ======================================================================
%
% SUBSECTION: DEFORMATION OF THE CLIFFORD COMPLETION
%
% ======================================================================

\subsection{Deformation theory of the Clifford completion}
% label removed: subsec:clifford-deformation
\index{Clifford completion!deformation theory}

The Clifford completion $G_g(\cA)$ depends on the central charge
parameter through $u = \kappa(\cA) \cdot \omega_g$. The variation of
$G_g$ with $\kappa$ is controlled by the deformation theory of
Clifford algebras.

\begin{proposition}[Deformation of the Clifford algebra;
\ClaimStatusProvedHere]
% label removed: prop:g9-clifford-deformation
\index{Clifford algebra!deformation theory}
\begin{enumerate}[label=\textup{(\roman*)}]
\item The family $\{\mathrm{Cl}_{2g}(u)\}_{u \in \Bbbk}$ is a flat
deformation of $\bigwedge^*(\Bbbk^{2g})$ over $\Bbbk$, with fiber
$\bigwedge^*(\Bbbk^{2g})$ at $u = 0$ and
$\mathrm{Mat}_{2^g}(\Bbbk)$ at $u \neq 0$.

\item The Hochschild cohomology controlling the deformation is
\begin{equation}% label removed: eq:g9-HH-clifford
\HH^2(\mathrm{Cl}_{2g}(u), \mathrm{Cl}_{2g}(u))
=
\begin{cases}
\Bbbk & \text{if } u = 0, \\
0 & \text{if } u \neq 0.
\end{cases}
\end{equation}
The one-dimensional space at $u = 0$ is generated by the
anticommutator deformation $\alpha_i\beta_j + \beta_j\alpha_i
\mapsto \delta_{ij}\,\epsilon$.

\item The deformation is unobstructed: $\HH^3(\bigwedge^*(\Bbbk^{2g}),
\bigwedge^*(\Bbbk^{2g})) = 0$ for the relevant component.
\end{enumerate}
\end{proposition}

\begin{proof}
\emph{(i)} The Clifford algebra is a filtered deformation of the
exterior algebra: the associated graded of $\mathrm{Cl}_{2g}(u)$
(filtered by word length) is $\bigwedge^*(\Bbbk^{2g})$ for all $u$,
including $u = 0$. Flatness is clear: $\mathrm{Cl}_{2g}(u)$ is free
of rank $2^{2g}$ as a $\Bbbk$-module for all $u$.

\emph{(ii)} For $u \neq 0$, $\mathrm{Cl}_{2g}(u) \cong
\mathrm{Mat}_{2^g}(\Bbbk)$, and $\HH^*(\mathrm{Mat}_n, \mathrm{Mat}_n)
\cong \HH^*(\Bbbk, \Bbbk) = \Bbbk$ concentrated in degree $0$ by
Morita invariance. In particular, $\HH^2 = 0$: matrix algebras are
rigid. For $u = 0$, $\HH^2(\bigwedge^*(\Bbbk^{2g})) = \Bbbk$,
generated by the single parameter $u$.

\emph{(iii)} The obstruction space $\HH^3$ for the exterior algebra
vanishes in the relevant weight: the deformation class $u$ has weight
$2$ (it changes the anticommutator from $0$ to $u \cdot \delta_{ij}$),
and the cup product $[u] \cup [u] \in \HH^4$ is the only potential
obstruction, which vanishes by the centrality of $u$.
\end{proof}

\begin{remark}[The wall-crossing at $u = 0$]
% label removed: rem:g9-wall-crossing
\index{wall-crossing!Clifford completion}
The deformation theory reveals that $u = 0$ is a \emph{wall}: for
$u \neq 0$ the algebra is rigid (a matrix algebra), while at $u = 0$
it acquires a one-dimensional deformation parameter. The Clifford
completion ``jumps'' across this wall, from a matrix algebra to an
exterior algebra, and the jump is irreversible (the exterior algebra
is not Morita equivalent to any matrix algebra). This is the algebraic
wall-crossing underlying the critical string dichotomy: the passage
through $\kappa = 0$ is a genuine phase transition in the derived
category.
\end{remark}

% ======================================================================
%
% SUBSECTION: GENUS-g TRACE FORMULA
%
% ======================================================================

\subsection{The genus-$g$ trace formula}
% label removed: subsec:trace-formula
\index{trace formula!Clifford completion}

In the topological regime, the Morita equivalence allows computation of
traces on $G_g$ via traces on $B_\Theta$.

\begin{proposition}[Trace formula; \ClaimStatusProvedHere]
% label removed: prop:g9-trace-formula
\index{trace formula!genus-g@genus-$g$}
When $u$ is invertible, any trace $\mathrm{tr}\colon B_\Theta \to \Bbbk$
extends to a trace on $G_g$ by the formula
\begin{equation}% label removed: eq:g9-trace-formula
\mathrm{Tr}_{G_g}(b \otimes C) =
2^g \cdot \mathrm{tr}(b) \cdot \mathrm{tr}_{\mathrm{Cl}}(C),
\end{equation}
where $\mathrm{tr}_{\mathrm{Cl}}\colon \mathrm{Cl}_{2g}(u)
\to \Bbbk$ is the reduced Clifford trace (normalized so that
$\mathrm{tr}_{\mathrm{Cl}}(1) = 1$).
\end{proposition}

\begin{proof}
Under the Morita equivalence
$G_g[u^{-1}] \cong \mathrm{Mat}_{2^g}(B_\Theta[u^{-1}])$, a trace on
$\mathrm{Mat}_{2^g}(R)$ is $\sum_{i=1}^{2^g}\mathrm{tr}(R_{ii})$.
The factor $2^g$ is the matrix size. The Clifford trace is the
composition of the matrix trace with the identification
$\mathrm{Cl}_{2g}(u) \cong \mathrm{Mat}_{2^g}$.
\end{proof}

\begin{corollary}[Genus-$g$ partition function via trace;
\ClaimStatusProvedHere]
% label removed: cor:g9-partition-trace
\index{partition function!trace formula}
For a Koszul chiral algebra $\cA$ with $\kappa(\cA) \neq 0$, the
genus-$g$ free energy satisfies
\begin{equation}% label removed: eq:g9-Fg-trace
F_g(\cA) = \frac{1}{2^g}\,\mathrm{Tr}_{G_g}\bigl(
(-1)^F \cdot q^{L_0}\bigr)
= \kappa(\cA) \cdot \lambda_g^{\mathrm{FP}},
\end{equation}
where the second equality uses the genus universality theorem
(Theorem~\ref{V1-thm:genus-universality}). The factor $1/2^g$ accounts
for the Morita multiplicity.
\end{corollary}

\begin{remark}[The partition function at $u = 0$]
% label removed: rem:g9-partition-u0
\index{partition function!gravitational regime}
In the gravitational regime ($u = 0$), the trace formula degenerates:
the Clifford trace is replaced by the supertrace on
$\bigwedge^*(\Bbbk^{2g})$, which is
$\mathrm{str} = \sum_{k=0}^{2g}(-1)^k\binom{2g}{k} =
(1-1)^{2g} = 0$ for $g \geq 1$. This vanishing is consistent with
$F_g(\mathrm{Vir}_0) = 0$: the gravitational partition function is
zero because the supertrace over the exterior algebra vanishes.
\end{remark}

% ======================================================================
%
% SUBSECTION: OPEN QUESTIONS
%
% ======================================================================

\subsection{Further directions}
% label removed: subsec:further-directions

\begin{remark}[Open questions]
% label removed: rem:g9-open-questions
\index{critical string!open questions}
\begin{enumerate}[label=\textup{(\alph*)}]
\item Does $G_g(\cA)$ carry a natural BV algebra structure
extending the BV/bar dictionary of Remark~\ref{V1-rem:bv-bar-bridge}? The quantum master
equation on $G_g$ should encode genus-$g$ sewing relations.

\item Is there a super-Clifford version for $\mathcal{N} = 1$ at
$c = 15$, replacing $\mathrm{Cl}_{2g}$ with a super-Clifford algebra
$\mathrm{sCl}_{2g|2g}$?

\item The Hodge--Clifford compatibility
(Proposition~\ref{V1-prop:hodge-clifford}) suggests a connection to the
Mukai Hodge structure on $K(\Sigma_g)$. Does the spectral sequence of
Construction~\ref{V1-constr:clifford-spectral} compute the Mukai pairing?

\item At $c = 13$, does the Clifford completion carry an enhanced
symmetry reflecting self-duality, beyond the $\Z/2$ of the Koszul
involution?

\item The wall-crossing at $u = 0$
(Remark~\ref{V1-rem:g9-wall-crossing}) should have a categorical lift:
a derived equivalence between the gravitational and topological
regimes that degenerates at $u = 0$. Can this be made precise using
the formalism of perverse sheaves on the $\kappa$-line?

\item The localization sequence
(Theorem~\ref{V1-thm:localization-sequence}) suggests a
``semi-classical limit'' interpretation: the gravitational regime is
the $u \to 0$ limit of the topological regime, analogous to the
$\hbar \to 0$ limit in deformation quantization. The secondary
anomaly $u = \kappa \cdot \omega_g$ plays the role of Planck's
constant.
\end{enumerate}
\end{remark}

%=======================================================================
% c = 26 no-ghost theorem as a chain-level Clifford-degeneration computation
%=======================================================================

\subsection{The $c=26$ BRST anomaly and Clifford degeneration}
\label{sec:critical-string-clifford-degeneration}

The preceding sections exhibited the $u\to 0$ degeneration of the
Clifford completion as the algebraic dichotomy associated to the
critical value $c=26$. The BRST matter--ghost system adds the
Goddard--Thorn filtration; the Clifford degeneration is the anomaly
calculation on the $E_1$ page, and the no-ghost theorem is the
collapse of that filtered complex.

\begin{theorem}[$c=26$ BRST anomaly cancellation and Clifford degeneration;
\ClaimStatusProvedHere]
\label{thm:c26-no-ghost-clifford-degeneration}
\index{no-ghost theorem!Clifford degeneration}
\index{critical string!chain-level no-ghost}
Let $\cA_{\mathrm{str}}$ be the bosonic string vertex algebra at
$c = 26$, assembled as the matter--ghost pair
$V_{\mathrm{matter}}^{c=26} \otimes V_{bc}^{c=-26}$ with total
central charge zero. Let $Q_{\mathrm{BRST}}$ be the BRST differential
and $\cH_{\mathrm{phys}}(\cA_{\mathrm{str}}) :=
H^\bullet(Q_{\mathrm{BRST}})$ the physical state space. Assume the
standard Goddard--Thorn hypotheses: the matter module has intercept one,
positive-energy grading, and nondegenerate transverse oscillator form,
and the Goddard--Thorn filtration has the Frenkel--Garland--Zuckerman
collapse
\[
E_1 \cong \cF_{\perp}\otimes H^\bullet(\Bbbk b_0\oplus \Bbbk c_0,
Q_0),
\qquad
E_1 = E_\infty .
\]
Here $Q_0$ is the zero-mode component of $Q_{\mathrm{BRST}}$ on the
$b_0,c_0$ sector.
Then:
\begin{enumerate}[label=\textup{(\roman*)}]
\item At $c^{\mathrm{matter}}=26$ (the critical string value), the
Clifford completion
(Section~\ref{app:anomaly-completed-topological-holography})
of the matter--ghost BRST pair $\cA_{\mathrm{str}}$ degenerates:
the effective handle anomaly
$\kappa_{\mathrm{eff}}^{\mathrm{BRST}} =
\kappa(\text{matter},c=26)+\kappa(\text{ghost},c=-26)$ is zero, so
$u_{\mathrm{eff}}=\kappa_{\mathrm{eff}}^{\mathrm{BRST}}\omega_g=0$ and
the handle Clifford algebra becomes
$\bigwedge^*(\Bbbk^{2g})$.
\item The $b_0,c_0$ zero modes form the usual odd Heisenberg pair
$\{b_0,c_0\}=1$, $b_0^2=c_0^2=0$. This pair is the contracting
homotopy on the longitudinal--ghost factor of the filtered BRST
complex, while the exterior handle algebra records the vanished
secondary anomaly.
\item The physical state space is canonically identified with the
transverse oscillator cohomology:
\[
\cH_{\mathrm{phys}}(\cA_{\mathrm{str}})
\;\cong\;
\cF_{\perp}.
\]
Longitudinal and nonzero ghost oscillators lie in
$Q_{\mathrm{BRST}}$-exact pairs.
\end{enumerate}
\end{theorem}

\begin{proof}
Effective anomaly vanishing at $c=26$: matter contributes
$\kappa(\text{matter}) = c^{\mathrm{matter}}/2 = 13$ (or
$\kappa = c$ in the convention of
\S~\ref{subsec:tholog-coda}); the $bc$-ghost system contributes
$\kappa(\text{ghost}) = -13$. Total $\kappa_{\mathrm{eff}}=0$,
so the secondary anomaly $u_{\mathrm{eff}}$ vanishes. Substituting
$u=0$ in the handle Clifford relations gives the exterior algebra
$\bigwedge^*(\Bbbk^{2g})$, by the gravitational-regime theorem above.

The $bc$ zero modes are not the handle Clifford generators; they form
the standard odd Heisenberg pair $\{b_0,c_0\}=1$ inside the BRST
complex. The operator $b_0$ is the contracting homotopy for the
longitudinal--ghost summand in the Goddard--Thorn filtered complex.
The assumed Frenkel--Garland--Zuckerman collapse identifies the
abutment with the transverse oscillator factor $\cF_\perp$, proving
the displayed cohomology isomorphism.
\end{proof}

\begin{remark}[From Clifford relation to Goddard--Thorn]
\label{rem:clifford-to-no-ghost}
The standard proof of no-ghost (Goddard--Thorn~\cite{GoddardThornNoGhost},
Frenkel--Garland--Zuckerman~\cite{FrenkelGarlandZuckerman}) is a
delicate spectral sequence argument. The Clifford-degeneration
statement isolates the anomaly calculation inside that proof: at
$c=26$, the matter and ghost modular characteristics cancel, the
handle Clifford algebra degenerates to an exterior algebra, and the
Goddard--Thorn filtration collapses onto transverse oscillators.
\end{remark}

\begin{remark}[Celestial consequence]
\label{rem:c26-celestial}
Via the universal celestial holography theorem
(Theorem~\ref{thm:universal-celestial-holography} in the
celestial-holography-core chapter), the $c=26$
no-ghost theorem is the celestial statement that the physical
soft-graviton sector of the bosonic string is the transverse
$\mathfrak{w}_{1+\infty}$ subalgebra, with longitudinal and ghost
soft modes $Q_{\mathrm{BRST}}$-exact. The Clifford degeneration
is precisely the shadow-class reduction $\mathbf C \rightsquigarrow
\mathbf G$ at the critical locus.
\end{remark}

\section{K3 class-\texorpdfstring{$\mathcal S$}{S} boundary constant}
\label{sec:csd-supplement}

\begin{proposition}[Class-$\mathcal S$ $A_1$ anomaly at twenty-four
punctures; \ClaimStatusProvedHere]
\label{rem:csd-1}
Let $\Sigma_{0,24}=(\mathbb P^1,D)$ with $|D|=24$, as for the
discriminant divisor of a semistable elliptic K3 surface with
twenty-four Kodaira $I_1$ fibres. The class-$\mathcal S$ theory
$\mathcal T[A_1,\Sigma_{0,24}]$ with maximal regular punctures has
\[
n_v = 21\cdot 3 = 63,\qquad
n_h = 22\cdot 4 = 88.
\]
Its four-dimensional and Schur-VOA central charges are
\[
c_{4d}
\;=\;
\frac{2n_v+n_h}{12}
\;=\;
\frac{126+88}{12}
\;=\;
\frac{107}{6},
\qquad
c_{2d}
\;=\;
-12c_{4d}
\;=\;
-214.
\]
\end{proposition}

\begin{proof}
A genus-zero surface with $n$ punctures decomposes into $n-2$
trinions and $n-3$ tubes. For $A_1$, each trinion $T_2$ contributes
$(n_v,n_h)=(0,4)$ and each tube contributes three vector multiplets,
so at $n=24$ one obtains $(n_v,n_h)=(63,88)$. The
Shapere--Tachikawa anomaly formula gives
$c_{4d}=(2n_v+n_h)/12=107/6$, equivalently
$c_{4d}=(5n-13)/6$ at genus zero. The Beem--Lemos--Liendo--Peelaers--
Rastelli--van~Rees Schur functor gives $c_{2d}=-12c_{4d}$.
For $n=4$ the same formula gives $c_{4d}=7/6$, the
$\mathrm{SU}(2)$ $N_f=4$ sanity check. The number $-312$ and the
factorisation $-12\cdot 26$ do not occur in this class-$\mathcal S$
anomaly calculation. The divisor $D$ varies in
\(\operatorname{Conf}_{24}(\mathbb P^1)/\operatorname{PGL}_2\); an
$M_{24}$-equivariant fibre is additional level structure on special
fibres, while the central-charge computation is the family statement.
\end{proof}
